\documentclass[11pt,a4paper]{article}
\usepackage[T1]{fontenc}
\usepackage[utf8]{inputenc}
\usepackage[margin=27mm]{geometry}
\usepackage{amsmath,amssymb,amsthm,mathtools,mathrsfs}
\usepackage{booktabs,array,longtable}
\usepackage[expansion=false]{microtype}
\usepackage[colorlinks=true,linkcolor=blue,citecolor=blue,urlcolor=blue]{hyperref}
\allowdisplaybreaks[2]
\setlength{\emergencystretch}{3em}
\numberwithin{equation}{section}
\newtheorem{maintheorem}{Theorem}
\renewcommand{\themaintheorem}{\Alph{maintheorem}}
\newtheorem{theorem}{Theorem}[section]
\newtheorem{proposition}[theorem]{Proposition}
\newtheorem{lemma}[theorem]{Lemma}
\newtheorem{corollary}[theorem]{Corollary}
\theoremstyle{definition}
\newtheorem{inputthm}[theorem]{Classical input}
\theoremstyle{remark}
\newtheorem{remark}[theorem]{Remark}
\newcommand{\Q}{\mathbb Q}
\newcommand{\C}{\mathbb C}
\newcommand{\R}{\mathbb R}
\newcommand{\Z}{\mathbb Z}
\newcommand{\OO}{\mathcal O}
\newcommand{\N}{\mathrm N}
\newcommand{\ind}{\mathbf 1}
\newcommand{\dd}{\,d}
\newcommand{\Rea}{\operatorname{Re}}
\newcommand{\E}{\mathcal E}
\newcommand{\rad}{\operatorname{rad}}
\newcommand{\sfstar}{\sideset{}{^*}\sum}
\newcommand{\GC}{\Gamma_{\C}}
\newcommand{\wt}{\widehat}
\newcommand{\ve}{\breve e}
\newcommand{\tauD}{\tau_{\mathrm{div}}}
\hypersetup{pdftitle={\detokenize{A zero-free half-plane for Hecke L-functions of nonzero angular type over Q(sqrt(-3)) and for the elliptic curves y^2 = x^3 + D}},pdfauthor={Jizhou Guo}}
\title{A zero-free half-plane for Hecke $L$-functions of nonzero angular type over $\Q(\sqrt{-3})$ and for the elliptic curves $y^2=x^3+D$}
\author{Jizhou Guo\\\texttt{mitsuha2021b@gmail.com}}
\date{October 8, 2026}
\begin{document}
\maketitle
\begin{abstract}
We extend the cubic-theta method of OpenAI's paper \emph{The Quasi-Riemann Hypothesis} of October 5, 2026 to Hecke characters of nonzero angular type over $K=\Q(\sqrt{-3})$. Every such unitary character, of arbitrary finite conductor, has a zero-free half-plane $\Rea s>11/12$. The half-plane is independent of the conductor and angular type. Pure horizontal jets at the cusps produce the required higher angular modes, while the two Poisson transformations preserve a fixed angular coset. Deuring's theorem gives $L(E_D,s)\ne0$ for $\Rea s>17/12$ for every $E_D:y^2=x^3+D$, $D\in\Q^\times$, in arithmetic normalization. The applications include power-saving sector and CM distribution estimates, polylogarithmic bounds for the least good prime in a fixed interval of positive CM Sato--Tate measure, and zero-free half-planes for all symmetric powers. The symmetric-power statement and the zero-frequency terms also use the finite-order $7/8$ theorem of the companion paper of September 30, 2026. The proof uses classical cubic-theta, reciprocity, quadratic large-sieve and Hecke theory, together with explicitly specified finite arithmetic and smooth-kernel statements of the source paper; all modifications of the source paper's cubic-theta argument involving angular type are derived here.
\end{abstract}
\tableofcontents

\section{Introduction and statements}
Fix the embedding $K\subset\C$ with $\sqrt{-3}=i\sqrt3$. We use the ideal-side convention for infinity types: a unitary Hecke character $\eta$ of conductor $\mathfrak f$ has angular type $k$ and norm parameter $t_0$ if
\[
 \eta((z))=(z/|z|)^k\N(z)^{it_0}\quad(z\equiv1\pmod{\mathfrak f}),\qquad k\in\Z,\quad t_0\in\R.
\]
The degree-two Hecke $L$-function is normalized to have functional equation centered at $1/2$.

\begin{maintheorem}\label{thm:A}
Let $K=\Q(\sqrt{-3})$. Every unitary Hecke character $\eta$ of finite conductor and nonzero angular type has
\[
 L_K(s,\eta)\ne0\qquad(\Rea s>11/12).
\]
The same half-plane applies to every real norm parameter $t_0$. More generally, if $\psi=\eta\N^u$, $u\in\R$, then $L_K(s,\psi)\ne0$ for $\Rea s>u+11/12$. In particular, for algebraic infinity type $(a,b)$ with $a\ne b$, the boundary is $(a+b)/2+11/12$.
\end{maintheorem}

\begin{maintheorem}\label{thm:B}
For every $D\in\Q^\times$, let $E_D/\Q$ be the elliptic curve $y^2=x^3+D$. Its arithmetic $L$-function satisfies
\[
 L(E_D,s)\ne0\qquad(\Rea s>17/12).
\]
The functional equation has center $s=1$. The nontrivial zeros, which lie in $1/2\le\Rea s\le3/2$, are confined to
\[
 7/12\le\Rea s\le17/12.
\]
\end{maintheorem}

The bounds in the intermediate mean squares are for a fixed character: their constants and thresholds may depend on its conductor and angular type. The location of the zero-free half-plane is common to all characters. Section~\ref{sec:applications} obtains estimates uniform in growing angular frequency from this common location and the analytic conductor, without imposing uniformity on those intermediate constants.

The curves $E_D$, with $D$ taken modulo rational sixth powers, form the family of all sextic twists of $y^2=x^3+1$, equivalently all elliptic curves over $\Q$ with $j=0$.

Hecke developed the analytic theory of Gr\"ossencharacter $L$-functions and the associated angular distribution of primes \cite{Hecke}. After removing a norm phase, the classical de la Vall\'ee-Poussin region has width $c$ divided by the logarithm of the analytic conductor; for a character of nonzero angular type the exceptional real quadratic zero cannot occur. A precise form is recorded in \cite[Lemma 2.1]{CPS}. Coleman's Vinogradov-type region \cite[Theorem 2]{Coleman}, in the formulation of \cite[Theorem 2.5]{Khale}, is
\begin{equation}\label{eq:classicalregion}
 \sigma\ge1-\frac{c_K}{\max\{\log\N\mathfrak f,(\log V)^{2/3}(\log\log V)^{1/3}\}},
 \qquad V=\max\{|k|,|t|\},
\end{equation}
for sufficiently large $V$, with the possible exception confined to $k=0$ and real finite character. For fixed nonzero type these widths tend to zero as the height increases, so these classical regions give no fixed half-plane inside the critical strip. Cossaboom's short-interval consequences of $\mathrm{GL}_2$ zero-density estimates use Coleman's region in the dihedral case \cite[Corollaries 1.3--1.4 and Section 9]{Cossaboom}.

The October 5, 2026 paper \cite{QRH12}, denoted \emph{p2} below, states $11/12$ for finite-order characters. The companion release states $7/8$ for finite-order Hecke characters of $K$ and for Dirichlet characters \cite[Theorem 1.1]{QRH78}. Other recent papers concern finite-order characters: \cite[Theorem 1.2]{PrimitiveRoots} states a width $10^{-6}$ over cyclotomic fields containing $\mu_{12}$, and \cite{Followup}, a draft in the public GitHub repository \nolinkurl{Rogerhu12/openai-math-extensions}, states a width $25/258$ over arbitrary number fields conditional on its theta, reflection and large-sieve inputs. These four statements concern finite-order characters. A second draft in that repository, dated 8 October 2026 \cite{FollowupSevenEighths}, states the half-plane $\Rea s>7/8$ for every unitary Hecke character, of finite or infinite order, over every number field, under a source hypothesis retaining the theta construction of \cite{PrimitiveRoots} and the estimates and moment arguments of \cite{QRH78}; its argument passes from a field $F$ to $F(\zeta_{12})$ and works with the finite-order twists of one fixed character. The present paper treats $K=\Q(\sqrt{-3})$ directly from p2: the angular modes are produced by pure jets of the cubic theta function, and the modifications of p2's argument involving the angular type are written out in full.

Our argument retains p2's M\"obius--sextic mean-square strategy. A Gauss--Jacobi conversion supplies one factor $\overline{\alpha}$, so a target of type $+k$ requires the theta mode $-(k+1)$. The finite-order case of p2 already has mode $-1$; the extension raises its derivative order to $m=k+1$. A pure Wirtinger jet realizes that mode at every cusp. Its reflection replaces the sextic row character by a quadratic one. The Goldmakher--Louvel large sieve acts on that finite quadratic family; the angular factor remains a row-independent column coefficient. Two Poisson transformations and cube inversion keep the same angular coset. Sixth powers of primes then extract the exponent $11/12$.

The proof uses the classical inputs stated in Section~\ref{sec:inputs}. The statements of p2 listed there with their hypotheses are used as printed; all modifications of its argument involving angular type, including reflection, transfer, descent and initialization, are derived here. The $7/8$ statement of the companion supplies the finite-order factors and zero-frequency terms in the applications. Section~\ref{sec:applications} proves power-saving sector and CM Sato--Tate estimates, a polylogarithmic least-prime bound uniform in the curve conductor, and zero-free half-planes for symmetric powers.

\section{Conventions and precise inputs}\label{sec:inputs}
Set $\omega=e^{2\pi i/3}$, $\OO=\Z[\omega]$, $\lambda=1+2\omega=i\sqrt3$, $\N(z)=z\bar z$ and $\alpha(z)=z/|z|$. The six units of $\OO$ map bijectively onto $(\OO/3\OO)^\times$. Since $\OO$ has class number one, every ideal prime to $3$ has a unique generator $n\equiv1\pmod3$, called primary; these generators multiply. Ideal indices below are represented by these generators. A star means squarefree. Row and Fourier-frequency indices range over elements of $\OO$, with all six associates when appropriate.

Fix a finite set $S$ of prime ideals containing those over $2,3$ and the target conductor. All primal ideal indices are prime to $S$ unless stated otherwise. In particular, the dual theta indices are \emph{not} subject to this restriction. Define
\[
 e(z)=\exp\bigl(2\pi i\operatorname{Tr}_{K/\Q}(z/\lambda)\bigr),\qquad
 \ve(z)=\exp\bigl(2\pi i(z+\bar z)\bigr)=e(\lambda z).
\]
For primary $n$ prime to $6$, put $\chi_n(x)=(x/n)_6$. Local powers are taken on the residue-field unit group and then extended by zero: $\chi_p^0(x)=\ind_{p\nmid x}$. This convention applies to the \emph{whole index}, including cube factors. Let $\mu$ be the ideal M\"obius function. For squarefree primary $n$ prime to $6$, set
\[
 \gamma_j(n)=\N(n)^{-1/2}\sum_{x\bmod n}\chi_n(x)^j e(x/n),\qquad
 G(n)=\overline{\chi_n(4)}\gamma_3(n).
\]
On other ray classes, $G$ means the finite extension in Input~\ref{inp:gauss}, rather than an imprimitive Gauss sum. In Gauss sums involving general primitive characters, the modulus is their primitive modulus. The empty modulus has $\gamma(\mathbf1)=1$ and $\mathbf1(0)=1$. For a fixed compact interval $I\subset(0,\infty)$, $\|W\|_{C^J(I)}$ is the maximum of the suprema of its ordinary derivatives of order at most $J$; for boxes we use total order.

\begin{inputthm}[Reciprocity and finite Gauss identities]\label{inp:gauss}
For primary elements prime to $S$, sextic reciprocity has a symmetric, $\{\pm1\}$-valued bicharacter $\mathcal R$ and a function $G$ with $|G|=1$ on a fixed finite ray group, with
\[
 \chi_b(a)=\mathcal R(a,b)\chi_a(b),\qquad
 G(ab)=G(a)G(b)\mathcal R(a,b)\quad((a,b)=1).
\]
For squarefree $n$ prime to $S$,
\begin{align}\label{eq:gauss}
 \gamma_2(n)^3&=\mu(n)\alpha(n),&
 \gamma_1(n)\gamma_2(n)&=\mu(n)\alpha(n)G(n),\\
 \gamma_1(n)\gamma_{-1}(n)&=\chi_n(-1),&
 \mu(n)\gamma_{-1}(n)&=\chi_n(-1)G(n)^{-1}\overline{\alpha(n)}\gamma_2(n).\label{eq:convert}
\end{align}
The function $G$ extends to the same finite ray group with its displayed cocycle, and
\[
 \chi_a(-1)\overline{G(a)}G(b)\mathcal R(a,b)=G(ba^{-1}).
\]
The quotient is taken in that ray group.
\end{inputthm}
These are the finite-symbol statements of p2 Lemma 4.1 and Appendix A.1, used as printed. Their hypotheses concern finite residue characters. The classical cubic reciprocity input is $(a/b)_3=(b/a)_3$ for coprime primary $a,b$, with the unit and $\lambda$ supplementary laws; an element congruent to $1\pmod9$ has trivial supplementary values. See \cite{Kubota} and \cite[(1.4)--(1.5), Section 2]{DR}. In particular $(\lambda/n)_3$ is a finite character supported at $3$. No angular character is put into a finite Fourier expansion.

\begin{inputthm}[Kubota--Patterson cubic theta]\label{inp:theta}
There is a smooth function $\theta$ on $\mathbb H^3=\C\times\R_{>0}$, with
\[
 \theta(gw)=\kappa(g)\theta(w),\quad
 \kappa\begin{pmatrix}a&b\\c&d\end{pmatrix}=(c/a)_3\quad(g\equiv I\pmod3,\ c\ne0),
\]
where $\kappa$ extends to $\Gamma_2=\langle\mathrm{SL}_2(\Z),\Gamma_1(3)\rangle$ and is $1$ on $\mathrm{SL}_2(\Z)$; it is $1$ when $c=0$. In particular $\theta$ is invariant under translations by $\Z+3\OO$.
Put $\gamma_0=I$, $\gamma_-={\left(\begin{smallmatrix}1&0\\\omega^2&1\end{smallmatrix}\right)}$, and $\gamma_+={\left(\begin{smallmatrix}1&0\\\omega&1\end{smallmatrix}\right)}$. Then
\begin{equation}\label{eq:thetaexp}
 \overline{\theta(\gamma_\sigma(z,v))}
 =\frac{3^{5/2}}2\ind_{\sigma=0}v^{2/3}
 +\sum_{0\ne\ell\in\lambda^{-4}\OO}d_\sigma(\ell)vK_{1/3}(4\pi|\ell|v)\ve(\ell z).
\end{equation}
Writing $t_\sigma$ for the corresponding unconjugated coefficients, the exact cusp identification is
\[
 t_0(\ell)=\tau(\ell),\quad
 t_-(\ell)=\omega^2\tau_1(\omega^2\ell)\ve(\ell),\quad
 t_+(\ell)=\omega\tau_2(\omega\ell)\ve(\ell),\quad
 d_\sigma(\ell)=\overline{t_\sigma(-\ell)}.
\]
Here $\tau,\tau_1,\tau_2$ are Patterson's coefficient sequences, given explicitly in \cite[(5.7), (5.13), (5.14)]{DR}. The coefficient $d_\sigma(\ell)$ vanishes unless
\[
 \ell=u\lambda^hnb^3,\quad u\in\OO^\times,\ h\ge-4,\quad n,b\text{ primary},\quad n\text{ squarefree},
\]
and in this representation
\begin{equation}\label{eq:d_bound}
 |d_\sigma(\ell)|\le27\,3^{h/6}\sqrt{\N(b)}.
\end{equation}
For squarefree primary $n$ and primary $b$ with $(nb,6)=1$, the coefficient of $\overline\theta$ at $\lambda^{-3}nb^3$ is exactly
\begin{equation}\label{eq:sourcecoeff}
 c_\theta(nb^3)=3^{5/2}\sqrt{\N(b)}\gamma_2(n)\chi_n(\lambda)^{-2}.
\end{equation}
All these Fourier series and their horizontal derivatives converge normally on compact subsets with $v>0$.
\end{inputthm}
This is Patterson's Theorem 8.1 and Tables II--III \cite{Patterson}, with the normalization and cusp representatives of \cite[(5.4)--(5.15), Appendix A]{DR}. These are unconditional coefficient and automorphy statements; the prime-Gauss-sum asymptotic of \cite{DR} is not used. We have stated precisely the coefficient identities and bounds needed here. For example, the two families in $t_0$ have magnitudes $3^{j/2+2}|b|$ when $h=3j-4$, $j\ge1$, and $3^{j/2+5/2}|b|$ when $h=3j-3$, $j\ge0$; the extra cusp coefficients at $h=-4$ have magnitude at most $9|b|$. These imply \eqref{eq:d_bound}.

\begin{inputthm}[Goldmakher--Louvel]\label{inp:GL}
Let $\mathfrak c$ be a fixed ideal of a number field $F$ and let $\{\rho_{\mathfrak a}\}$ be a quadratic Hecke family indexed by squarefree ideals prime to $\mathfrak c$. Precisely, each member is a primitive character of trivial infinity type and order dividing $2$; there exist a finite group $G_0$, an ideal homomorphism $[\cdot]$ into $G_0$ and $C:G_0^2\to\{\pm1\}$ such that
\[
 \rho_{\mathfrak a}(\mathfrak b)=\rho_{\mathfrak b}(\mathfrak a)C([\mathfrak a],[\mathfrak b])\quad((\mathfrak a,\mathfrak b)=1);
\]
and, for coprime $\mathfrak a,\mathfrak b$ of the same class, the primitive character inducing $\rho_{\mathfrak a}\rho_{\mathfrak b}$ has conductor exactly $\mathfrak a\mathfrak b$. Then for $M,N\ge1$ and arbitrary complex coefficients $\beta(\mathfrak b)$,
\[
 \sfstar_{\N\mathfrak a\le M}\left|\sfstar_{\N\mathfrak b\le N}\beta(\mathfrak b)\rho_{\mathfrak b}(\mathfrak a)\right|^2
 \ll_{F,\mathfrak c,\epsilon}(MN)^\epsilon(M+N)\sfstar_{\N\mathfrak b\le N}|\beta(\mathfrak b)|^2.
\]
\end{inputthm}
This is \cite[Definition 1, Theorem 1.1]{GL}. Section~\ref{sec:completed} verifies the three properties for the family actually used and the direction of this inequality.

\begin{inputthm}[Hecke theory and the prime ideal theorem]\label{inp:Hecke}
A primitive unitary Hecke character of $K$ of type $(k,t_0)$ with $k\ne0$ has an entire $L$-function. With $q=3\N\mathfrak f$,
\[
 \Lambda(s,\eta)=q^{s/2}\GC(s+|k|/2-it_0)L_K(s,\eta),\qquad
 \GC(s)=2(2\pi)^{-s}\Gamma(s),
\]
satisfies $\Lambda(s,\eta)=\epsilon_\eta\Lambda(1-s,\bar\eta)$, $|\epsilon_\eta|=1$, up to an $s$-independent normalization. For $k=0$ the only pole occurs for a pure norm character, at $s=1+it_0$. The completed functions have order one after poles are removed. The Euler product is absolutely convergent for $\Rea s>1$. Also, for any fixed finite $S$,
\[
 \#\{\mathfrak p\notin S:Y/2<\N\mathfrak p\le Y\}\sim\frac{Y}{2\log Y},
\]
so it is at least $Y/(4\log Y)$ for $Y$ sufficiently large.
\end{inputthm}
See \cite{Hecke,Tate}, \cite[Section 5.3]{Watkins}, and the prime ideal theorem in \cite[Chapter 5]{IK}. The fixed-field prime ideal theorem supplies the prime rows used here.

The remaining printed inputs of p2 are Lemma 4.2 (primitive finite-character row Poisson with a radial Schwartz weight, positive row scale and an arbitrary extra ideal exclusion), Lemma 6.1 (finite periodic masks of the scalar theta coefficient sequence), and Appendix B, Lemmas B.1--B.3 (smooth Mellin separation, the row-dependent kernel bound and weighted recombination for arbitrary finite coefficient arrays). We give the row Poisson formula in \eqref{eq:rowpoisson}; the smooth-kernel hypotheses and proofs are included in Appendix~\ref{app:smooth}. The arithmetic part of Lemma 6.4 is the support and bound \eqref{eq:d_bound}. Its transformed-weight assertion is replaced by our higher-mode calculation. The angular counterparts of p2 Propositions 4.4, 5.1, 5.2, 5.4 and 6.2 are proved below, together with the finite cusp construction underlying Lemma 6.3.

\section{Angular characters and pure cusp jets}\label{sec:jet}
\begin{lemma}[Primary factorization and norm shifts]\label{lem:factor}
A unitary character of type $(k,t_0)$ may be written, on primary ideals prime to $3\mathfrak f$,
\[
 \eta(n)=\nu(n)\alpha(n)^k\N(n)^{it_0},
\]
where $\nu$ is a finite ray character modulo $3\mathfrak f$. Moreover
\[
 L_K(s,\bar\eta)=\overline{L_K(\bar s,\eta)},\qquad
 L_K(s,\eta\N^z)=L_K(s-z,\eta).
\]
\end{lemma}
\begin{proof}
The primary generators multiply, so after removal of $\alpha^k\N^{it_0}$ the quotient is a multiplicative ideal character. It is $1$ on generators congruent to $1\pmod{3\mathfrak f}$, hence factors through the finite ray class group. The displayed identities follow by absolute convergence of Euler series in a right half-plane and then analytic continuation. Continuous characters of $\C^\times\simeq\R_{>0}\times S^1$ are $\alpha^k\N^z$, with $k\in\Z$ and $z\in\C$: this follows by taking logarithms on the positive factor and using integer Fourier frequencies on the circle. For unitary characters $z=it_0$. Removing $\N^{it_0}$ preserves the real parts of zeros. For algebraic type $(a,b)$, $z=(a+b)/2$ and $k=a-b$. We may therefore prove the mean square for $\nu\alpha^k$ with $k>0$; negative types follow by conjugation.
\end{proof}

Fix henceforth $k\ge1$, set $\tau=\alpha^{-k}$, and choose a finite ray group $T$, supported at $S$, through which all finite target and reciprocity factors pass, including the finitely many $S$-supported Kummer residue characters with exponents read modulo six. Recursion will use the finite coset
\[
 \mathscr C_k=\{\tau\chi:\chi\in\widehat T\}.
\]
It need not be a group. For $\xi\in\mathscr C_k$ put
\begin{equation}\label{eq:axi}
 a_\xi(n)=\overline{\alpha(n)}\gamma_2(n)\xi(n).
\end{equation}
The completed coefficient has angular mode $-m$, where $m=k+1\ge2$.

\begin{lemma}[Pure Wirtinger jet at a cusp]\label{lem:jet}
Let $a\delta'-bc=1$, $c\ne0$. Suppose smooth functions satisfy, near $z=0$ for every $v>0$,
\[
 F(z+a/c,v)=\zeta F_H(z',v'),\qquad
 z'=-\delta'/c-\frac{\bar z}{c^2(v^2+z\bar z)},\quad
 v'=\frac v{\N(c)(v^2+z\bar z)}.
\]
Here $\zeta$ is constant and $|\zeta|=1$. For every integer $m\ge1$,
\begin{equation}\label{eq:jet}
 \partial_{\bar z}^mF(z+a/c,v)\big|_{z=0}
 =\zeta(-c^{-2}v^{-2})^m
 (\partial_{z'}^mF_H)(-\delta'/c,1/(\N(c)v)).
\end{equation}
There are no height derivatives or mixed terms in this identity.
\end{lemma}
\begin{proof}
Use independent formal variables $z,\bar z$ for the real smooth Taylor jet of order $m$. Modulo the ideal $(z)$,
\[
 z'=-\delta'/c-\bar z/(c^2v^2),\quad
 \bar z'=-\bar\delta'/\bar c,\quad v'=1/(\N(c)v).
\]
Only $z'$ varies, and it varies affinely. Thus the coefficient of $\bar z^m$ in the composite Taylor jet comes solely from $\partial_{z'}^mF_H/m!$, with factor $(-c^{-2}v^{-2})^m$. A smooth Taylor remainder $o(|z|^m)$ does not change this coefficient. This proves \eqref{eq:jet}. In particular, any cusp term depending on height alone has zero pure horizontal jet for $m\ge1$. At $m=0$ it survives.
\end{proof}

There is a representation-theoretic description of this calculation. The complexified real Lie algebra of $\mathrm{SL}_2(\C)$ is $\mathfrak{sl}_2\oplus\mathfrak{sl}_2$. Let $E=e_1$ and $\bar E=e_2$ be the two horizontal raising directions. For the Weyl inversion $w$, $\operatorname{Ad}(w)\bar E=-f_2=K_0-E$, where $K_0=e_1-f_2\in\mathfrak k_\C$ for $\mathfrak k=\mathfrak{su}_2$. The two summands commute, so $[K_0,E]=0$. A right $\mathrm{SU}_2$-invariant spherical vector is killed by $R(K_0)$; hence $R((K_0-E)^m)=(-1)^mR(E^m)$ on it. In upper half-space this vector is $v^m\partial_{\bar z}^m\theta$, with horizontal rotation weight $m$. The commutation explains why no height or mixed derivatives appear. The cusp reflection uses precisely this pointwise jet identity.

\section{Mellin reflection in an arbitrary positive mode}\label{sec:mellin}
For an expansion \eqref{eq:thetaexp}, put
\[
 J_m^-(s;F)=\int_0^\infty\partial_{\bar z}^mF(0,v)v^{2s+m-2}\dd v,\qquad
 G_m(s)=\Gamma(s+m/2-1/6)\Gamma(s+m/2+1/6).
\]
A plus superscript means $\partial_z^m$ instead.
\begin{lemma}[Mellin normalization]\label{lem:mellin}
Termwise integration in an absolutely convergent range gives
\begin{equation}\label{eq:modeMellin}
 J_m^-(s;F)=\frac{i^m}{4(2\pi)^{2s}}G_m(s)
 \sum_{\ell\ne0}d(\ell)\alpha(\ell)^{-m}\N(\ell)^{-s}.
\end{equation}
For $J_m^+$ replace $\alpha^{-m}$ by $\alpha^m$.
For a finite periodic mask $F(A)$, zero-extended at the whole index, define
\[
 \mathcal T_m(s;F)=\sum_{n\ \mathrm{sf},b}\gamma_2(n)F(nb^3)\alpha(nb^3)^{-m}
 \N(n)^{-s}\N(b)^{-3s+1/2}.
\]
The scalar theta mask realization of p2 Lemma 6.1 then has
\begin{equation}\label{eq:maskedMellin}
 J_m^-(s;\Theta_F)=\frac{3^{5/2}}4
 \left(\frac{27}{(2\pi)^2}\right)^sG_m(s)\mathcal T_m(s;F).
\end{equation}
\end{lemma}
\begin{proof}
The Fourier derivative is $(2\pi i\bar\ell)^m$, while the Bessel Mellin integral \cite[(10.43.19)]{DLMF} is
\[
 \int_0^\infty v^{2s+m-1}K_{1/3}(4\pi|\ell|v)\dd v
 =2^{2s+m-2}(4\pi|\ell|)^{-2s-m}G_m(s),
 \quad\Rea s>1/6-m/2.
\]
Their product is \eqref{eq:modeMellin}. For the full finite mask $F$ in this lemma put
\[
 \phi(A)=\ind_{A\equiv1(3),(A,S)=1}\chi_A(\lambda)^2F(A),
\]
extended by zero outside the indicated set. Finite Fourier inversion realizes the scalar masked theta series as a finite sum of translates. In Section~\ref{sec:cusps}, for $F=F_0\prod_{p\in\mathcal P}\chi_p^{j_p}$, the fixed factor and the good-prime local masks are Fourier-expanded separately; their product is the same full mask. The pure jet supplies $\alpha^{-m}$. The condition $(\lambda^3\ell,S)=1$ forces $h=-3$ in the support decomposition, and $\lambda^3\ell\equiv1\pmod3$ forces the unit to be $1$. Thus only $\ell=\lambda^{-3}nb^3$ remains. The factor $\chi_{nb^3}(\lambda)^2=\chi_n(\lambda)^2$ cancels the coefficient factor in \eqref{eq:sourcecoeff}, which proves the finite mask realization explicitly. Since $\alpha(\lambda^{-3})=i$, the product $i^m\alpha(\lambda^{-3})^{-m}$ is $1$, yielding \eqref{eq:maskedMellin}. All $n,b$ sums are absolutely integrable for $\Rea s>1$: the two norm exponents are $\Rea s>1$ and $3\Rea s-1/2>1$. This justifies Fubini, including shared primes between $n$ and $b$.
\end{proof}

Applying \eqref{eq:jet} and changing $v=(\N(c)u)^{-1}$ gives the exact cusp identity
\begin{equation}\label{eq:MellinFE}
 J_m^-(s;F(\cdot+a/c))=(-1)^m\zeta\alpha(c)^{-2m}\N(c)^{1-2s}
 J_m^+(1-s;F_H(\cdot-\delta'/c)).
\end{equation}
Indeed $c^{-2m}\N(c)^{m+1-2s}=\alpha(c)^{-2m}\N(c)^{1-2s}$. The norm scale is independent of $m$.

For $V\in C_c^\infty(I)$ define $\wt V(t)=\int_0^\infty V(x)x^t\dd x/x$ and
\begin{align}\label{eq:Rm}
 R_m(t)&=\frac{\Gamma((m+1)/2+t-1/6)\Gamma((m+1)/2+t+1/6)}
 {\Gamma((m+1)/2-t-1/6)\Gamma((m+1)/2-t+1/6)},\\
 V_m^\sharp(x)&=\frac1{2\pi i}\int_{(0)}\wt V(-t)R_m(t)
 \left(\frac{(2\pi)^4x}{27}\right)^{-t}\dd t.\label{eq:kernel}
\end{align}
\begin{lemma}[Contour and transformed weight]\label{lem:contour}
For every fixed mask and $m\ge1$, $\mathcal T_m(s;F)$ is entire and has polynomial vertical growth on every fixed strip. Reflection of its smoothed completion gives, for each cusp shift, the scalar
\begin{equation}\label{eq:scalar}
 \frac{(-i)^m}{81}\zeta\alpha(c)^{-2m}
\end{equation}
multiplying
\[
 \sum_{\ell\ne0}\frac{d(\ell)\alpha(\ell)^m}{\sqrt{\N(\ell)}}
 \ve(-\delta'\ell/c)V_m^\sharp(\N(\ell)X/\N(c)^2).
\]
For $A>0$, $j\ge0$, there is $J=\lceil4A\rceil+j+6$ such that
\begin{equation}\label{eq:kernelbound}
 |(x\partial_x)^jV_m^\sharp(x)|\ll_{m,A,j,I}\|V\|_{C^J(I)}\min(x^{1/4},x^{-A}).
\end{equation}
\end{lemma}
\begin{proof}
Fix the finite shifts and cusp matrices. At infinity their differentiated Fourier series have exponential decay. At zero \eqref{eq:jet} expresses them as a power of $v$ times a differentiated Fourier series at height $1/(\N(c)v)$; every constant term is killed. Thus the Mellin integrand has bounds $O(v^Be^{-cv})$ at infinity and $O(v^{-B}e^{-c'/v})$ at zero, with constants allowed to depend on this fixed mask. The Mellin integral is entire and bounded on any fixed strip. Dividing by $G_m$ in \eqref{eq:maskedMellin} gives at most exponential vertical growth by Stirling.

On $\Rea s=1+\delta$, the source series is absolutely bounded. On $\Rea s=-\delta$, \eqref{eq:MellinFE} and the absolutely convergent dual series, together with the gamma quotient, give polynomial growth. Choose $B_0>2+\delta$ and an integer $M_0$ exceeding both boundary exponents. The function $f(s)=\mathcal T_m(s)/(s+B_0)^{M_0}$ is holomorphic in the intervening strip, bounded on its vertical boundaries and at most exponential inside. Apply the maximum principle on rectangles to $f(s)e^{\epsilon s^2}$: its horizontal-edge bound tends to zero as the rectangle height tends to infinity, while its vertical-edge bound is $C e^{\epsilon C_\delta}$. Let the height tend to infinity and then $\epsilon\downarrow0$. This proves polynomial growth throughout the strip and justifies shifting a compact smooth Mellin weight without residues.

Start from
\[
 \frac1{2\pi i}\int_{(\sigma>1)}\wt V(s-1/2)\mathcal T_m(s;F)X^{s-1/2}\dd s.
\]
Use \eqref{eq:maskedMellin} and \eqref{eq:MellinFE}, then put $t=1/2-s$. The gamma ratio is exactly \eqref{eq:Rm}; the norm factors give $\N(\ell)X/\N(c)^2$. The remaining constants are
$(-1)^mi^m3^{-5/2}27^{-1/2}=(-i)^m/81$, proving \eqref{eq:scalar}. The dual series is first expanded where $\Rea t>1/2$; each resulting kernel is then moved to $\Rea t=0$.

The nearest numerator pole of $R_m$ is $-(m+1)/2+1/6\le-5/6$, and reciprocal gamma factors are entire. Shift to $-1/4$ for $x\le1$ and to $A$ for $x\ge1$. Stirling gives $|R_m(\sigma+it)|\ll_{m,A}(1+|t|)^{4\sigma}$ on these lines. Mellin integration by parts $J$ times absorbs this growth and the factor $|t|^j$ from Euler derivatives, proving \eqref{eq:kernelbound}. Contour movement is performed separately for each fixed mask. The subsequent moving-row estimates require only \eqref{eq:kernelbound}, whose constants depend on $m$ and the indicated test parameters, not on a moving mask or row.
\end{proof}

\section{Finite masks and all cusp classes}\label{sec:cusps}
For the fixed factor $F_0$ in the masks below put
\[
 \phi_0(A)=\ind_{A\equiv1(3),(A,S)=1}\chi_A(\lambda)^2F_0(A),
\]
extended by zero outside the indicated set. Choose $L$ once, with support in $S$, divisible by all fixed conductor and exclusion primes in $S$ and high enough powers at $2,3$ to make $\phi_0$ periodic. It may be enlarged to be divisible by $18$. Put $M=\lambda^{12}L^4$. Neither $L$ nor $M$ changes with the good-prime masks. The normalized finite Fourier coefficient is
\[
 \wt\phi_0(h)=\N(L)^{-1}\sum_{A\bmod L}\phi_0(A)e(-hA/L),
 \qquad |\wt\phi_0(h)|\le\|F_0\|_\infty\le1.
\]

\begin{lemma}[Cusp reduction and multiplier]\label{lem:cusps}
Let $\mathcal P$ be a finite set of good primary primes, $j_p\in\{0,\ldots,5\}$, and
\[
 F(A)=F_0(A)\prod_{p\in\mathcal P}\chi_p^{j_p}(A).
\]
Here $F_0$ is a finite periodic mask with a fixed period supported on $S$, zero-extended at its fixed exclusions, with $|F_0|\le1$; this includes a finite ray character multiplied by a class function of modulus at most one. The finite theta mask is a sum of translates at
\[
 z_{\boldsymbol h}=\lambda^2\left(h_0/L+\sum_{p\in\mathcal P}h_p/p\right),\quad
 h_0\bmod L,\ h_p\bmod p,
\]
with coefficients $\wt\phi_0(h_0)\prod_pC_{p,j_p}(h_p)$, where
\begin{equation}\label{eq:finiteFourier}
 C_{p,j}(h)=\begin{cases}
 q^{-1/2}\chi_p(-1)^j\gamma_j(p)\chi_p(h)^{-j},&j\ne0,\ h\ne0,\\
 0,&j\ne0,\ h=0,\\
 -q^{-1},&j=0,\ h\ne0,\\
 1-q^{-1},&j=0,\ h=0,
 \end{cases}\qquad q=\N(p).
\end{equation}
For the active set $\mathcal A=\{p:h_p\ne0\}$, set $r=\prod_{p\in\mathcal A}p$. The denominator ideal of $z_{\boldsymbol h}$ is the product of $(r)$ and the denominator ideal of $\lambda^2h_0/L$, even when $(h_0,L)\ne1$. Use active representatives $h_p\equiv0\pmod{M^2}$, as constructed below. Choose $z_{\boldsymbol h}=a/c$ in lowest terms with $a\equiv1\pmod3$ if $\lambda\mid c$, and $c\equiv1\pmod3$ otherwise, and define $c_0=c/r$. Throughout the multiplier formulas and the definitions of $\sigma_p,\epsilon_p,D_0$ below, $c_0$ is this normalized element, so $c=c_0r$ exactly. For fixed $h_0$ and fixed $r\bmod M^2$, there are a cusp coefficient sequence $d\in\{d_0,d_-,d_+\}$, an additive character $\psi$ and a unit $\kappa_0$, all independent of the active values $h_p$, such that, for $x=\lambda^4\ell$,
\begin{align}\label{eq:cuspCRT}
 \kappa(g_1)&=\kappa_0\prod_{p\mid r}\chi_p(a)^2,\\
 \ve(-\delta'\ell/c)&=\psi(x)\prod_{p\mid r}e(\epsilon_ph_p^{-1}x/p),\qquad
 \epsilon_p=-\bigl((\lambda^3c/p)\sigma_p\bigr)^{-1},\quad \sigma_p=\lambda^2c/p.\notag
\end{align}
The normalization unit in $c$ depends only on $(h_0,r\bmod3)$. Since the active primes are primary, $r\equiv1\pmod3$, so here this unit depends only on $h_0$.
\end{lemma}
\begin{proof}
Define $C_{p,j}(h)=q^{-1}\sum_{x\bmod p}\chi_p^j(x)e(-hx/p)$. Character orthogonality at $h=0$, and $x=h^{-1}y$ at $h\ne0$, prove \eqref{eq:finiteFourier}. Finite Fourier inversion of $\phi_0$ and these masks gives the stated translates. Their constant terms sum to zero, because $\phi_0(0)=0$. This construction never Fourier-expands $\alpha^{-m}$.

At an active good prime $p$, the summand $\lambda^2h_p/p$ has valuation $-1$, while every other summand is integral at $p$; it cannot cancel. At a prime in $S$, all good-prime summands are integral, and the negative valuation, if present, is exactly that of $\lambda^2h_0/L$. This proves the denominator assertion without any coprimality assumption on $h_0$.

By CRT choose each active representative $h_p\equiv0\pmod{M^2}$. Changing a representative only translates by $3\OO$, a period of $\theta$. Normalize $a/c$ as in the statement and set $c_0=c/r$. To identify the normalization unit, start with any fixed generator $c_{00}$ of the denominator ideal of $\lambda^2h_0/L$ and the preliminary denominator $c_{00}r$. Since
\[
 a=\lambda^2c\left(h_0/L+\sum_{p\mid r}h_p/p\right),
\]
its fixed high residues at $S$ are determined by $h_0,r\bmod M^2$. In particular the preliminary numerator residue selecting the normalization unit is $(\lambda^2c_{00}h_0/L)r\pmod3$ in the numerator-normalized case; the good-prime terms vanish in that congruence. Thus this unit depends only on $(h_0,r\bmod3)$.

Choose $\delta'$ by CRT with
\[
 \begin{cases}
 a\delta'\equiv1\pmod{q^{v_q(Mc_0)}}&q\mid c_0,\\
 \delta'\equiv0\pmod{q^{v_q(M)}}&q\mid M,\ q\nmid c_0,\\
 a\delta'\equiv1\pmod r.&
 \end{cases}
\]
Set $b=(a\delta'-1)/c$ and $g={\left(\begin{smallmatrix}a&b\\c&\delta'\end{smallmatrix}\right)}$. Then $b\equiv0\pmod{q^{v_q(M)}}$ for $q\mid c_0$ and $b\equiv-c^{-1}$ at the other $S$ primes. These prescriptions include the inert prime over $2$ and every target conductor prime: division by $c_0$ is compensated by the modulus $Mc_0$. There is no assumption that the $2$-part is squarefree or that $h_0$ is a unit.

For the three possible classes choose
\begin{equation}\label{eq:Hcases}
 H=\begin{cases}
 I,&3\mid c,\\
 {\left(\begin{smallmatrix}1&0\\u&1\end{smallmatrix}\right)},\quad u=\pm\lambda,\ u\equiv c\pmod3,&v_\lambda(c)=1,\\
 {\left(\begin{smallmatrix}u&-1\\1&0\end{smallmatrix}\right)},\quad u\equiv a\pmod3,&(c,\lambda)=1,
 \end{cases}
\end{equation}
where in the last case $u$ comes from a fixed residue list. The congruences just imposed show directly that $g_1=gH^{-1}\equiv I\pmod3$. Invariance under $\mathrm{SL}_2(\Z)$ and translations by $\Z+3\OO$ identifies the \emph{function} $\theta(Hw)$ with one of $\theta(\gamma_\sigma w)$. More explicitly, lower translations are conjugates of upper translations by the integral inversion matrix. The lower parameters $\lambda$ and $-\lambda$ differ from $\omega^2$ and $\omega$, respectively, by elements of $\Z+3\OO$. In the last case, applying the integral inversion to $T_u w$ reduces to lower parameter $-u$, whose class modulo $\Z+3\OO$ is $0,\omega$ or $-\omega$. This identifies all three cusp functions with their multipliers.

Here is the multiplier calculation in each case, with cubic symbols interpreted using the supplementary laws when a numerator is not primary:
\begin{equation}\label{eq:multipliercases}
 \kappa(g_1)=\begin{cases}
 (c_0/a)_3(a/r)_3,&3\mid c,\\
 (-u/(a-ub))_3((c_0/u)/a)_3(a/r)_3,&v_\lambda(c)=1,\\
 (a/c_0)_3(a/r)_3,&(c,\lambda)=1.
 \end{cases}
\end{equation}
In the first case this is reciprocity applied to $(c/a)_3$. In the second let $A=a-ub$; the first column of $g_1$ is $(A,c-u\delta')$. The determinant relation gives
\[
 a(c-u\delta')\equiv-u\pmod A,\qquad bc\equiv-1\pmod a.
\]
Since $a,A$ are primary, reciprocity gives $(a/A)_3=(A/a)_3=(-ub/a)_3$. Combining these two congruences gives the second line. In the third case the first column is $(-b,-\delta')$ and
\[
 \kappa(g_1)=(\delta'/(-b))_3,\quad \delta'\equiv a^{-1}\pmod b,\quad
 -bc=1-a\delta'\equiv1\pmod9.
\]
To justify reciprocity if $\lambda\mid a$, write $a=\varepsilon\lambda^v a_0$ with $a_0$ primary and prime to $\lambda$. The last congruence makes the supplementary values of $\varepsilon$ and $\lambda$ trivial, and the determinant gives $-bc\equiv1\pmod{a_0}$. Thus $(a/(-bc))_3=1$. It follows that $\kappa(g_1)=(a/c)_3$, the third line. The degenerate $a=0$ case has unit $c$ and is checked directly using the integral inversion and translations. A zero $b$ or a zero lower-left entry of $g_1$ is treated by the defining value $\kappa=1$, so every symbol used has a nonzero denominator.

Every factor other than $(a/r)_3$ in \eqref{eq:multipliercases} is fixed by the prescribed $S$-adic residues. Factor such a numerator as a unit, a power of $\lambda$, and primary primes $q$ in $S$. Reciprocity reduces each $(q/a)_3$ to the residue of $a$ modulo $q$; the unit and $\lambda$ supplementary values depend only on $a\bmod9$. In the second case the same statement applies to $A=a-ub$, whose residue modulo $9$ is fixed. In the third case $(a/c_0)_3$ depends only on $a\bmod\rad c_0$. These moduli divide $M$, since $c_0$ divides $L$ up to a unit. This proves the first line of \eqref{eq:cuspCRT} and its ray-class consistency, including the prime above $2$.

Finally, $\ve(-\delta'\ell/c)=e(-\delta'x/(\lambda^3c))$. Put $D_0=\lambda^3c_0$. CRT at denominator $D_0r$ gives the fixed $S$ term
\[
 \psi(x)=e(-\delta'_0r^{-1}x/D_0),\quad \delta'_0=\delta'\bmod D_0.
\]
At $p\mid r$, $a\equiv\sigma_ph_p$ and $\delta'_p=(\sigma_ph_p)^{-1}$; its additive term is exactly the second line of \eqref{eq:cuspCRT}. All inverses exist. The character $\psi$, the cusp function and $c_0$ therefore come from a fixed finite list once the indicated ray class is fixed.
\end{proof}

\begin{proposition}[Full-mask higher-mode reflection]\label{prop:reflection}
For the masks $F$ of Lemma~\ref{lem:cusps}, in particular with $|F_0|\le1$, and for $m\ge1$, $V\in C_c^\infty((0,\infty))$, set
\[
 T_m(X;F,V)=\sum_{n\ \mathrm{sf},b}\frac{\gamma_2(n)\alpha(nb^3)^{-m}F(nb^3)}{\sqrt{\N(n)}\N(b)}
 V(\N(n)\N(b)^3/X).
\]
This finite sum equals the sum, over $h_0\bmod L$ and
$\{p:j_p\ne0\}\subseteq\mathcal A\subseteq\mathcal P$, of
\begin{equation}\label{eq:reflection}
 C_{m,h_0,\mathcal A}\sum_{\ell\ne0}
 \frac{d(\ell)\alpha(\ell)^m\psi(\lambda^4\ell)}{\sqrt{\N(\ell)}}
 \prod_{p\in\mathcal A}B_{p,j_p}(\lambda^4\ell)
 V_m^\sharp(\N(\ell)X/\N(c)^2),
\end{equation}
where
\begin{equation}\label{eq:Bpj}
 B_{p,j}(x)=\begin{cases}
 \chi_p(x)^{-j-2},&j\notin\{0,4\},\\
 q^{-1/2}(-1+q\ind_{p\mid x}),&j=4,\\
 q^{-1/2}\chi_p(x)^{-2},&j=0,
 \end{cases}
\end{equation}
and the exact scalar is
\begin{equation}\label{eq:Cscalar}
 C_{m,h_0,\mathcal A}=\frac{(-i)^m}{81}\alpha(c)^{-2m}\wt\phi_0(h_0)\bar\kappa_0
 \prod_{p\in\mathcal P\setminus\mathcal A}(1-q_p^{-1})
 \prod_{p\in\mathcal A}\chi_p(\sigma_p)^{-2}\omega_{p,j_p},\quad |C_{m,h_0,\mathcal A}|\le1/81,
\end{equation}
with
\[
 \omega_{p,j}=\begin{cases}
 \chi_p(-1)^j\gamma_j(p)\gamma_{j+2}(p)\chi_p(\epsilon_p)^{-j-2},&j\notin\{0,4\},\\
 \gamma_4(p),&j=4,\\
 -\gamma_2(p)\chi_p(\epsilon_p)^{-2},&j=0\text{ active}.
 \end{cases}
\]
All $\omega_{p,j}$ have modulus one. Both $n$ and $b$ may contain the same prime.
\end{proposition}
\begin{proof}
Apply Lemmas~\ref{lem:jet}--\ref{lem:contour} to the finite translate sum of Lemma~\ref{lem:cusps}. The denominator, including its normalization unit, is independent of active $h_p$, so the new factor $\alpha(c)^{-2m}$ can be taken outside the local sums. For an active prime put $y=h_p^{-1}$. The product of the Fourier character and conjugate multiplier becomes
\[
 \chi_p(h_p)^{-j}\chi_p(a)^{-2}=\chi_p(\sigma_p)^{-2}\chi_p(y)^{j+2}.
\]
If $j\ne0,4$, its nontrivial Gauss sum gives $\omega_{p,j}\chi_p(x)^{-j-2}$. If $j=4$, the remaining character is principal on nonzero residues and the additive sum is $-1+q\ind_{p\mid x}$. If $j=0$ is active, its Fourier coefficient is $-q^{-1}$, giving $q^{-1/2}\chi_p(x)^{-2}$. These computations also hold at $x=0$. An inactive prime can only have $j=0$ and supplies $1-q^{-1}$. Combining them with \eqref{eq:scalar} gives \eqref{eq:reflection}--\eqref{eq:Cscalar}. Absolute convergence of the dual sum follows from \eqref{eq:d_bound} and \eqref{eq:kernelbound}.
\end{proof}

When the moving squarefree row $k_0$ has exponent $j_p=1$ at its primes, every such prime is active and
\[
 B_{p,1}(\lambda^4\ell)=\chi_p(\lambda^4\ell)^3.
\]
Fix a modulus $\mathfrak m_*$ supported on $S$, divisible by $(M^2)$, $24\OO$ and the modulus of the ray group in Input~\ref{inp:gauss}. Partition the moving squarefree rows into classes modulo $\mathfrak m_*$ and apply Lemma~\ref{lem:GLfamily} class by class. Write the active product as $r=r_{\mathrm{fixed}}k_0$, where $r_{\mathrm{fixed}}$ is the product of the other active prime generators, and define $c_*=c/k_0=c_0r_{\mathrm{fixed}}$. This element is fixed for $h_0$ and the fixed active part. The angular phase separates exactly as $\alpha(c_*k_0)^{-2m}=\alpha(c_*)^{-2m}\alpha(k_0)^{-2m}$ within each class. The second factor is row-only; $\alpha(\ell)^m$ is column-only. This separation is the interface to the quadratic large sieve.

\section{The completed mean square}\label{sec:completed}
\begin{lemma}[The quadratic family and its three properties]\label{lem:GLfamily}
For squarefree primary $k_0$ prime to $S$, let $e_{k_0}=0,1$ according as $\N(k_0)\equiv1,3\pmod4$, and let $\kappa_\lambda$ be the nontrivial character of $(\OO/\lambda)^\times$. The character on elements
\[
 \rho_{k_0}(x)=(x/k_0)_2\kappa_\lambda(x)^{e_{k_0}}
\]
induces a primitive quadratic Hecke character of trivial infinity type and conductor $k_0\lambda^{e_{k_0}}$. When $e_{k_0}=0$ the $\lambda$ factor is omitted. This family satisfies all three conditions of Input~\ref{inp:GL}. Consequently, for arbitrary squarefree-column coefficients supported on $\N(n)\asymp U$,
\begin{equation}\label{eq:GLactual}
 \sfstar_{\N(k_0)\le H_0}\left|\sfstar_n\beta(n)\chi_{k_0}(n)^3\right|^2
 \ll_{S,\delta}(H_0U)^\delta(H_0+U)\sum_n|\beta(n)|^2.
\end{equation}
The columns may have a fixed $S$-part.
\end{lemma}
\begin{proof}
The quadratic symbol of $\omega$ is $1$, since $\omega$ is a square in $\mu_6$. Its value at $-1$ is $(-1)^{e_{k_0}}$. The corresponding values of $\kappa_\lambda$ are $1$ and $-1$. Their product is therefore trivial on all six units and defines an ideal character of trivial infinity type and order dividing $2$: this verifies condition (1), including the ideal rather than element interpretation. Each prime of $k_0$ has a nontrivial residue-field quadratic character, and $\kappa_\lambda$ is nontrivial precisely when $e_{k_0}=1$; CRT shows the conductor is exactly the one claimed.

Partition the primary generators into classes modulo the fixed $\mathfrak m_*$ specified above. On primary arguments the $\lambda$ factor is $1$, so Input~\ref{inp:gauss} gives
\[
 \rho_{k_0}(n)=\chi_{k_0}(n)^3=\mathcal R(k_0,n)\rho_n(k_0).
\]
Use a fixed common refinement with the reciprocity ray group if needed. The primary-generator residue map is multiplicative and has finite image; taking $C=\mathcal R$ proves condition (2) with the required finite group homomorphism. This partition fixes all unit, $2$-adic and $\lambda$ supplementary factors.

For coprime rows $k_1,k_2$ in the same class, $e_{k_1}=e_{k_2}$, so their $\lambda$ factors cancel on taking the primitive inducer. The product remains nontrivial at every prime of $k_1k_2$ and has no other conductor prime. Its primitive conductor is $k_1k_2$, verifying condition (3). This also applies if the class has $e=0$: there is no artificial zero-extension at $\lambda$ in its primitive conductor.

For primary columns $n$, $\kappa_\lambda(n)=1$ and $\rho_{k_0}(n)=\chi_{k_0}(n)^3$. Theorem 1.1 of \cite{GL} has the opposite row/column notation; on each pair of classes reciprocity changes it by one fixed sign, or equivalently one can use Hilbert-space duality. Collisions give zero on both sides and do not affect this argument. The finite partition proves \eqref{eq:GLactual}. If $n$ has primes in $S$, fix their product first (at most $2^{|S|}$ possibilities), and move its quadratic value to the row scalar. The remaining ideal column is prime to the fixed family modulus. Norm ranges change only by fixed factors.
\end{proof}

For $\xi=\alpha^{-k}\chi\in\mathscr C_k$, define, with whole-index zero extensions,
\begin{align}\label{eq:Psi}
 \Psi_{r,f}(A)&=\xi(A)\chi_A(r)\chi_A(f)^4,\\
 T(X;r,f)&=\sum_{n\ \mathrm{sf},b}
 \frac{\overline{\alpha(nb^3)}\gamma_2(n)\Psi_{r,f}(nb^3)}{\sqrt{\N(n)}\N(b)}
 V_*(\N(n)\N(b)^3/X),\quad V_*(y)=\sqrt yW(y).\label{eq:Tcompleted}
\end{align}
\begin{proposition}[Completed energy]\label{prop:completed}
Fix $k\ge1$, $T,S$, $C_0\ge1$, $\epsilon>0$ and $I\Subset(0,\infty)$. For $D\ge2$, $1\le H,X,\N(f)\le D^{C_0}$, with $f$ squarefree primary and prime to $S$,
\begin{equation}\label{eq:completedbound}
 \sum_{0<\N(r)\le H}|T(X;r,f)|^2
 \ll_{k,T,S,C_0,\epsilon,I}D^\epsilon\|W\|_{C^{32}(I)}^2
 \left(H+\frac{H^2\N(f)}X\right).
\end{equation}
\end{proposition}
\begin{proof}
We give the squarefree-row argument with an arbitrary nonzero auxiliary element $g$ in place of $f$. Choose primary generators for prime ideals outside $3$, $\lambda$ above $3$, and extend multiplicatively. A squarefree element row is $u_0tk_0$, where $t$ collects primes in $gS$, and $k_0$ is squarefree and prime to $gS$. Fix $t,u_0$ and set $H_0=H/\N(t)\ge1$; empty ranges are discarded. After sextic reciprocity, the finite local mask has exponent
\[
 j_p\equiv v_p(u_0tk_0)+4v_p(g)\pmod6.
\]
At $p\mid k_0$, this is $1$. Other exponents are fixed on $tg$; exponent zero still means the nonunit exclusion. The remaining $S$ factor comes from a fixed finite list. The angular mode is $m=k+1$ throughout, so Proposition~\ref{prop:reflection} applies.

Within each fixed class of $k_0$ modulo $\mathfrak m_*$, write $\ell=u\lambda^hnb^3$, $h\ge-4$. The coefficients satisfy \eqref{eq:d_bound}; $n,b$ need not be prime to $S$. The new column coefficient before fixed local factors is
\[
 A_0(n,b)=\frac{d(u\lambda^hnb^3)\alpha(u\lambda^hnb^3)^m\psi(u\lambda^{h+4}nb^3)}{3^{h/6}\sqrt{\N(b)}},
 \qquad |A_0|\le27.
\]
It is independent of $k_0$. Remove row scalars $\alpha(c_*k_0)^{-2m}$ and $\chi_{k_0}(u\lambda^{h+4})^3$ in the $\ell^2$ norm. Each branch is reduced to
\begin{equation}\label{eq:Sbranch}
 \mathscr S_h(k_0)=\sum_{n\ \mathrm{sf},b}
 \frac{A(n,b)\chi_{k_0}(nb)^3}{3^{h/3}\sqrt{\N(n)}\N(b)}
 V_m^\sharp\left(\frac{3^h\N(n)\N(b)^3H_0^2}{Y\N(k_0)^2}\right),
\end{equation}
initially with $|A|\le27a$, $a=1$, $Y=\N(c_0)^2H_0^2/X$.

For a fixed active $p\nmid k_0$, the conductor scale incurs $q^2$, $q=\N(p)$. The exceptional local factor splits exactly as
\[
 B_{p,4}(u\lambda^{h+4}nb^3)=-q^{-1/2}+q^{1/2}\ind_{p\mid n}+q^{1/2}\ind_{p\nmid n,\ p\mid b}.
\]
Reindex $n=pn'$ in the second term and $b=pb'$ in the third, moving $\chi_{k_0}(p)^3$ to a row scalar. The complete amplitude/scale budget is
\[
\begin{array}{c|c}
 \text{local term}&(a^2,Y)\text{ after insertion}\\\hline
 j\ne0,4&(a^2,q^2Y)\\
 j=0&(a^2/q,q^2Y)\\
 j=4:\ -q^{-1/2}&(a^2/q,q^2Y)\\
 j=4:\ p\mid n&(a^2,qY)\\
 j=4:\ p\nmid n,\ p\mid b&(a^2/q,Y/q)
\end{array}
\]
The second reindex keeps $p\nmid n'$ but permits $p\mid b$; the third keeps $p\nmid n$ but permits $p\mid b'$. Thus no shared-prime term is lost. Distinct-prime reindexes commute. At $p\mid t$ the exponent is odd and costs at most $q^2$ in $a^2Y$. At $p\mid g$, $p\nmid t$, $j=4v_p(g)$ is $0,4$ or $2$; the cost is at most $q$ in the first two cases and $q^2$ in the third, where $v_p(g)\ge2$. Hence every branch satisfies
\begin{equation}\label{eq:budget}
 a^2\le1,\qquad a^2Y\ll_{T,S}\frac{H_0^2\N(t)^2\N(g)}X.
\end{equation}
There are at most $C_{T,S}\tauD(tg)^3$ branches. All these divisor losses can be absorbed in preassigned small powers of $D$ in the polynomial parameter range; for example $Y\ll D^{6C_0}$ suffices.

For each fixed $b$, write $\chi_{k_0}(nb)^3=\chi_{k_0}(b)^3\chi_{k_0}(n)^3$. The first factor is row-only and has modulus at most one, including its collision zeros; absorb it into the majorized Mellin coefficient, not into the column sequence. Take dyadic ranges $\N(n)\asymp U$, $\N(b)\asymp B$, and $z_*=3^hUB^3/Y$. Mellin-separate only $\N(n)/U$ in the weight of \eqref{eq:Sbranch}. By \eqref{eq:kernelbound}, its row-dependent coefficients have common majorant
\[
 C_{m,I}\|W\|_{C^{32}}(1+z_*)^{-3}(1+|t|)^{-2}.
\]
The restriction $(k_0,g)=1$ is dropped by positivity when applying the large sieve; its fixed family modulus remains supported on $S$, and no moving-$g$ constant is introduced. For fixed $b,t$, the column sequence is $A(n,b)\N(n)^{-1/2}(\N(n)/U)^{it}$ and is independent of the row. Its square norm is $O(a^2)$ by ideal counting. Weighted Cauchy--Schwarz in the Mellin integral and in $b$, using $\sum_{\N(b)\asymp B}\N(b)^{-1}\ll1$, followed by \eqref{eq:GLactual}, gives
\begin{equation}\label{eq:blockenergy}
 \sum_{\N(k_0)\le H_0}^{*}|\mathscr S_{h;U,B}(k_0)|^2
 \ll a^2\|W\|_{C^{32}}^2 3^{-2h/3}(H_0U)^\delta(H_0+U)(1+z_*)^{-6}.
\end{equation}
A fixed $S$-part of $n$ is handled as in Lemma~\ref{lem:GLfamily}; $b$ is fixed before this application and requires no finite decomposition.

Since $h\ge-4$, $U\le81Yz_*$. Choose $\delta$ small relative to $\epsilon,C_0$; the square root of \eqref{eq:blockenergy} is bounded by
\[
 C aD^{\epsilon_1}\|W\|_{C^{32}}3^{-h/3}(\sqrt{H_0}+\sqrt U)(1+3^hUB^3/Y)^{-2}.
\]
For dyadic $U,B\ge1$ and every $Z>0$,
\[
 \sum_{U,B}(1+UB^3/Z)^{-2}\ll\log^2(2+Z),\qquad
 \sum_{U,B}\sqrt U(1+UB^3/Z)^{-2}\ll\sqrt Z.
\]
To prove the second inequality, sum $U$ on either side of $U=Z/B^3$ and then sum $B^{-3/2}$; if $Z<1$, the whole sum is $O(Z^2)\le O(\sqrt Z)$. For the first, count the dyads below $UB^3=Z$ and sum geometric tails outside. Summing $h\ge-4$ gives convergent factors $3^{-h/3}$ and $3^{-h/2}$. The $\ell^2$ triangle inequality therefore bounds each branch by
\[
 C aD^{2\epsilon_1}\|W\|_{C^{32}}(\sqrt{H_0}+\sqrt Y).
\]
Sum the branches and $t$ using \eqref{eq:budget} and divisor bounds. With losses allocated in advance this proves
\begin{equation}\label{eq:sfcompleted}
 \sum_{s\ \mathrm{sf},\ \N(s)\le H}|T(X;u_0s,g)|^2
 \ll D^\epsilon\|W\|_{C^{32}}^2(H+H^2\N(g)/X).
\end{equation}

For general rows write uniquely $r=u_0sv^2$ with $s$ squarefree; $s,v$ may share primes. Whole-index zero extensions give
\[
 \chi_n(u_0sv^2)\chi_n(f)^4=\chi_n(u_0s)\chi_n(fv^2)^4,
\]
since $8\equiv2\pmod6$, with both sides zero at collisions. Apply \eqref{eq:sfcompleted} with row range $H/\N(v)^2$ and $g=fv^2$, increasing the fixed polynomial range if necessary. Both terms acquire $\N(v)^{-2}$. The sum over $v$ converges, $\sum_v\N(v)^{-2}=\zeta_K(2)$; the six units contribute a fixed factor. This proves \eqref{eq:completedbound}.
\end{proof}

\section{Two Poisson transformations preserve the angular coset}\label{sec:transfer}
This section allows $\tau$ to be any completely multiplicative unit-modulus function on ideals prime to $S$. The theta input is not asserted for this generality. Define, for $\xi\in\tau\widehat T$,
\begin{equation}\label{eq:energy}
 \E(H,X,F;\xi,W)=\frac1{XF}\sfstar_{F\le\N(f)<2F}
 \sum_{0<\N(r)\le H}\left|\sfstar_n a_\xi(n)\chi_n(r)\chi_n(f)^4W(\N(n)/X)\right|^2.
\end{equation}
Let $\E_R$ add the column restriction $(n,R)=1$. Only primal ideal columns and labels are prime to $S$; rows and frequencies are not restricted at $S$.

\begin{lemma}[Paired conversions and exclusion removal]\label{lem:paired}
For coprime squarefree primary $a,b$, $a_\xi(ab)=a_\xi(a)a_\xi(b)\chi_b(a)^4$. If the finite functions have Fourier expansions on $T$
\[
 \chi(z)G(z)=\sum_{\chi_1}l_{\chi_1}\chi_1(z),\qquad
 \chi_1(z)G(z^{-1})=\sum_{\chi'}c_{\chi'}\chi'(z),
\]
then, for coprime paired columns prime to $S$,
\begin{align}\label{eq:pair1}
 a_{\tau\chi}(u_1)\overline{a_{\tau\chi}(u_2)}\gamma(\chi_{u_1}\bar\chi_{u_2})
 &=\mu(u_1)\mu(u_2)\sum_{\chi_1}l_{\chi_1}(\tau\chi_1)(u_1)\overline{(\tau\chi_1)(u_2)},\\
 \mu(z_1)\mu(z_2)(\tau\chi_1)(z_1)\overline{(\tau\chi_1)(z_2)}\gamma(\bar\chi_{z_1}\chi_{z_2})
 &=\sum_{\chi'}c_{\chi'}a_{\tau\chi'}(z_1)\overline{a_{\tau\chi'}(z_2)}.\label{eq:pair2}
\end{align}
The Fourier coefficient $\ell^1$ norms are at most $\sqrt{|T|}$. For $R_0=\prod_{p\mid R,\ p\notin S}p$,
\begin{equation}\label{eq:exclusion}
 \E_R(H,X,F;\xi,W)\le\tauD(R_0)\sum_{d\mid R_0}
 \E(H,X/\N(d),F\N(d);\xi,W).
\end{equation}
\end{lemma}
\begin{proof}
The CRT identities of Input~\ref{inp:gauss} give the multiplicative formula when $\tau=1$; multiply by $\tau(ab)=\tau(a)\tau(b)$. For coprime $a,b$, CRT gives $\gamma(\chi_a\bar\chi_b)=\mathcal R(a,b)\gamma_1(a)\gamma_{-1}(b)$ and the conjugate-character version with indices $-1,1$. Applying \eqref{eq:convert} and the quotient identity in Input~\ref{inp:gauss} then gives the finite factor $\chi(u_1u_2^{-1})G(u_1u_2^{-1})$ in the first direction and $\chi_1(z_1z_2^{-1})G(z_2z_1^{-1})$ in the second. Multiply both sides of each identity by $\tau(u_1)\bar\tau(u_2)$ or $\tau(z_1)\bar\tau(z_2)$ and expand only the displayed finite function. This proves \eqref{eq:pair1}--\eqref{eq:pair2}, with $\tau$, not $\bar\tau$, in both outputs. Parseval and Cauchy--Schwarz on $T$ give the coefficient norms. Every common factor cancels \emph{before} absolute values:
\begin{equation}\label{eq:commonphase}
 \tau(dv_1)\overline{\tau(dv_2)}=\tau(v_1)\overline{\tau(v_2)}.
\end{equation}
For exclusion removal, inclusion--exclusion and the multiplicative formula give the exact inner-sum identity
\[
 C_R(X;r,f)=\sum_{d\mid R_0,(d,f)=1}\mu(d)a_\xi(d)\chi_d(r)\chi_d(f)^4C_1(X/\N(d);r,df).
\]
Here $\chi_n(d)^4$ supplies the shared-prime exclusion; terms with $(d,f)\ne1$ vanish. Exterior coefficients have modulus at most one. Cauchy--Schwarz in $d$ and the injective map $f\mapsto df$ prove \eqref{eq:exclusion}; $XF$ is unchanged. This argument is valid for $H,X>0$, $F\ge1$.
\end{proof}

We use p2 Lemma 4.2 with its stated finite-character hypotheses: for a primitive finite residue character $\rho$ modulo $\mathfrak m$, radial Schwartz $\Phi$, $H>0$ and an arbitrary extra exclusion ideal $R$,
\begin{equation}\label{eq:rowpoisson}
 \sum_r\rho(r)\ind_{(r,R)=1}\Phi(\N(r)/H)
 =\frac{H\gamma(\rho)}{\sqrt{\N\mathfrak m}}
 \sum_{d\mid\rad R}\frac{\mu(d)\rho(d)}{\N(d)}
 \sum_h\bar\rho(h)\wt\Phi\left(\frac{H\N(h)}{\N(d)\N\mathfrak m}\right),
\end{equation}
where
\[
 \wt\Phi(|w|^2)=\frac2{\sqrt3}\int_\C\Phi(|z|^2)e(-zw)\dd x\dd y.
\]
The nonprincipal character has zero Fourier coefficient at $h=0$; for the primitive principal character the zero term is $H\wt\Phi(0)\prod_{p\mid R}(1-\N(p)^{-1})$. Formula \eqref{eq:rowpoisson} follows also by inclusion--exclusion and lattice Poisson on residue classes. In every use below, the row character is a finite sextic residue character, and $\tau$ is an external column coefficient.

Choose a nonnegative radial Schwartz $\Phi\ge1$ on $[0,1]$ with $\wt\Phi$ supported in $[0,C_\Phi]$. Such a weight is obtained by squaring a real radial function with smooth compactly supported Fourier transform, nonzero at the origin, and rescaling. Let $\mathcal A(W)$ be \eqref{eq:energy} with column scale $L$ and row cutoff replaced by $\Phi(\N(r)/H)$, including $r=0$.

\begin{proposition}[Transfer]\label{prop:transfer}
Fix $I=[u,v]\Subset(0,\infty)$, $j\ge0$, $\epsilon>0$ and $C_0\ge1$. If
\[
 D\ge2,\quad1\le H,L,F,\Sigma\le D^{C_0},\qquad\max(H,LF)\le\Sigma,
\]
then
\begin{equation}\label{eq:transfer}
 \mathcal A(W)\ll_{T,S,I,j,\epsilon,C_0,\Phi}
 D^\epsilon\Sigma\|W\|_{C^{4j+12}(I)}^2
 \left(1+\sup\frac{\E(H',X',F';\xi',U)}{\Sigma'}\right),
\end{equation}
where $\xi'\in\tau\widehat T$, $\Sigma'=X'F'$ and
\begin{equation}\label{eq:children}
 H',X',F'\ge1,\qquad H'\le\frac{HL}{\Sigma F},\quad
 H'/\Sigma'\le H/\Sigma,\quad\Sigma'\le L,
\end{equation}
and $U\in C_c^\infty(I')$, $I'=[u/16,4v]$, $\|U\|_{C^j}\le1$. An empty supremum is zero.
\end{proposition}
\begin{proof}
We record both transformations and the cancellation that makes their child kernels identical.

\emph{The first Poisson transformation.} Expand $\mathcal A$ and write $n_i=Cu_i$ with $C=(n_1,n_2)$. Apply \eqref{eq:rowpoisson} to the primitive finite character $\chi_{u_1}\bar\chi_{u_2}$ with exclusion $C$. Its zero frequency requires $u_1=u_2=1$ and contributes $O_I(H\|W\|_\infty^2)$. At nonzero frequency use \eqref{eq:pair1}. Remove $(u_1,u_2)=1$ by $\sum_{t\mid u_1,u_2}\mu(t)$ and write $u_i=tx_i$. Formula \eqref{eq:commonphase} cancels the $C,t$ angular factors exactly. The remaining column sum is
\begin{equation}\label{eq:firstP}
 P_{C,d,t}(y;U)=\sfstar_n\mu(n)(\tau\chi_1)(n)\ind_{(n,t)=1}
 \bar\chi_n(y)\chi_n(C)^4\chi_n(d)U(\N(n)/\ell),\quad \ell=L/(\N(C)\N(t)),
\end{equation}
where $d\mid C$ and the original labels and frequencies enter only through $y=hf^2$. Their ranges are
\[
 (f,Ct)=(h,t)=1,\quad
 0<\N(h)\le\frac{C_\Phi v^2L^2\N(d)}{H\N(C)^2}.
\]
Put
\[
 w_{C,d}=\frac{H\N(C)}{\N(d)L^2F},\quad
 Y_{C,d}=\frac{c_I\Sigma LF\N(d)}{H\N(C)^2},\quad
 c_I=\max(1,4C_\Phi v^2).
\]
Then $\N(y)\le Y_{C,d}$. The exact coupled kernel, in scaled column variables, is
\[
 (x_1x_2)^{-1/2}W(x_1)\bar W(x_2)
 \wt\Phi\left(\frac{H\N(h)\N(C)^2}{L^2\N(d)x_1x_2}\right).
\]
On its nonzero support, the Fourier argument is bounded by $C_\Phi$, so any prescribed finite set of derivatives is bounded uniformly in all indices. The map $(f,h)\mapsto hf^2$ has divisor-bounded multiplicity. Positivity, $\Phi\ge1$ and the smooth-kernel lemmas of Appendix~\ref{app:smooth} therefore reduce the nonzero contribution to the positive form
\begin{equation}\label{eq:Qform}
 Q_{\tau\chi_1}(U)=\sum_{C,t\ \mathrm{sf},(C,t)=1,\ \N(Ct)\le2vL}
 \sum_{d\mid C}w_{C,d}\sum_y\Phi(\N(y)/Y_{C,d})|P_{C,d,t}(y;U)|^2.
\end{equation}
More precisely, a bound $Q_{\tau\chi_1}(U)\le M\|U\|_{C^{j_0}(I_*)}^2$ for all $\chi_1$, with $I_*=[u/2,2v]$, implies
\begin{equation}\label{eq:firsttransfer}
 \mathcal A(W)\ll D^{\epsilon_0}(\Sigma+M)\|W\|_{C^{2j_0+4}(I)}^2.
\end{equation}
All the signs $\mu(d)\mu(t)$ from this first stage have been bounded before passing to the positive form; the common angular factors have already canceled.

\emph{The second Poisson transformation.} Fix $C,d,t$ in \eqref{eq:Qform}, put $Y=Y_{C,d}$, and apply \eqref{eq:rowpoisson} in $y$. Write the expanded columns as $gz_i$. Let $e\mid g$ come from the row exclusion and $w\mid(z_1,z_2)$ from coprimality inversion. Formula \eqref{eq:pair2} converts the paired coefficients back into the same coset. At nonzero frequency $h$, the new coefficient is
\begin{equation}\label{eq:asharp}
 a^\sharp(n)=a_{\tau\chi'}(n)\ind_{(n,tg/e)=1}\chi_n(deh)\chi_n(Cew)^4,
\end{equation}
with new column length $X'_0=\ell/\N(gw)$ and exterior factor
\[
 \frac{Y\mu(e)\mu(w)c_{\chi'}\N(g)}{\N(e)\ell}.
\]
The ranges are
\[
 (g,Ct)=(w,gCth)=1,\quad \N(gw)\le2v\ell,\quad
 0<\N(h)\le\frac{4C_\Phi v^2\ell^2\N(e)}{Y\N(g)^2}.
\]
The exact kernel is
\[
 U_0(x_1)\bar U_0(x_2)
 \wt\Phi\left(\frac{Y\N(h)\N(g)^2}{\N(e)\ell^2x_1x_2}\right),\quad U_0(x)=x^{-1/2}U(x).
\]
No angular factor remains in the Fourier test. In particular, cancellation of each common phase in \eqref{eq:commonphase} holds before any signed preimages are combined.

Regroup by
\[
 R_0=tg/e,\qquad f'=Cew,\qquad r'=deh.
\]
Both $R_0,f'$ are squarefree and coprime. Direct substitution gives
\begin{align}\label{eq:secondscales}
 X'_0&=L/(\N(R_0)\N(f')),\\
 w_{C,d}\frac{Y_{C,d}\N(g)}{\N(e)\ell}&=\frac{c_I\Sigma\N(R_0)}{L^2},\notag\\
 \frac{Y_{C,d}\N(h)\N(g)^2}{\N(e)\ell^2}
 &=\frac{c_I\Sigma F\N(r')\N(R_0)^2}{HL}.\notag
\end{align}
Thus the coefficient \eqref{eq:asharp}, the scale and the kernel depend only on $(R_0,f',r',\chi')$, not on their preimage. For fixed $R_0$, there are at most $\tauD(R_0)$ admissible $t\mid R_0$, namely those with $(t,r')=1$; the following sign sum is independent of $t$. At a prime $p\mid f'$ the three locations $p\mid C,e,w$ have signed total
\begin{equation}\label{eq:preimages}
 (1+\ind_{p\mid r'})-\ind_{p\mid r'}-\ind_{p\nmid r'}=\ind_{p\mid r'}.
\end{equation}
The first location permits $p\mid d$ exactly when $p\mid r'$; the second forces $p\mid r'$; the third requires $p\nmid h$, hence $p\nmid r'$. These are all preimages. Since their coefficients and kernels agree by \eqref{eq:secondscales}, their signs can be combined before taking absolute values. We are left with $f'\mid r'$. The finite coefficients $c_{\chi'}$ do not depend on a preimage, and no $\tau(C),\tau(e)$ or $\tau(w)$ survives.

\emph{Child ranges and norms.} Retain only $R_0$ for which the original kernel has a nonzero term, and keep this set when enlarging $f',r'$ ranges. In dyadic ranges $R\le\N(R_0)<2R$, $F'\le\N(f')<2F'$, set
\[
 X'=L/(RF'),\quad \Sigma'=L/R,\quad H'=HL/(\Sigma FR^2).
\]
Kernel support, with our choice of $c_I$, gives
\[
 \N(f')\le\N(r')\le\frac{HL}{\Sigma F\N(R_0)^2}.
\]
Every retained $R_0$ therefore has $H'\ge1$. For every $j_0\mid R_0$,
\begin{equation}\label{eq:childpositive}
 \frac{X'}{\N(j_0)}\ge\frac{\Sigma F\N(R_0)^2}{HR\N(j_0)}\ge\frac{\Sigma F}{H}\ge1.
\end{equation}
Exclusion removal \eqref{eq:exclusion} thus never creates a column scale below $1$ and preserves $H',\Sigma'$. The remaining inequalities are exactly \eqref{eq:children}. All child parameters are at most $D^{C_0}$. The ratio $X'_0/X'$ is in $[1/4,1]$; the common-test support is in $[u/8,2v]\Subset I'$, with uniformly bounded finite derivatives.

Let the supremum in \eqref{eq:transfer} be $S_j$. For each retained $R_0$, \eqref{eq:exclusion} and the definition of $S_j$ give a common-test squared-sum bound
\[
 \ll D^\delta(\Sigma')^2S_j\|U\|_{C^j(I')}^2.
\]
The row-dependent kernel lemma turns this into a $C^{2j+4}$ bound. There are $O(R)$ possible $R_0$ in the block, and the exterior coefficient is $\Sigma R/L^2$, so the net cost is
\[
 \frac{\Sigma R}{L^2}\,R\,(L/R)^2=\Sigma.
\]
The second zero frequency comes only from identical coprime columns. Its common angular factor is $|\tau(g)|^2=1$, and its total is bounded by
\[
 C\Sigma\|U\|_\infty^2
 \sum_C\frac{\tauD(C)}{\N(C)^2}
 \sum_{\N(t)\le2vL}\frac1{\N(t)}
 \ll_{I,\delta}\Sigma D^\delta\|U\|_\infty^2.
\]
The first sum converges and the second is logarithmic; no lower bound on $Y_{C,d}$ is needed. Assign sufficiently small norm exponents to all divisor bounds, absorb the finitely many dyadic losses, and take $\delta<\epsilon/20$. We obtain the $C^{2j+4}$ bound for $Q$. In \eqref{eq:firsttransfer} take $j_0=2j+4$; the resulting derivative order is $4j+12$, proving \eqref{eq:transfer}.
\end{proof}

This proves the transfer for a general unit-modulus multiplicative $\tau$ but does not give its completed estimate. Only $\tau=\alpha^{-k}$ is connected to Proposition~\ref{prop:completed} by the theta realization.

\section{Cube inversion and finite descent}\label{sec:descent}
Restore $\tau=\alpha^{-k}$, $k\ge1$. Write $\Sigma=XF$.
\begin{lemma}[Cube inversion and reduction]\label{lem:cube}
For $H\le\Sigma$ in a fixed polynomial parameter range,
\begin{equation}\label{eq:cubeidentity}
 X^{-1/2}\sfstar_n a_\xi(n)\chi_n(r)\chi_n(f)^4W(\N(n)/X)
 =\sum_h\frac{\mu(h)\bar\alpha(h)^3\Psi_{r,f}(h)^3}{\N(h)}T(X/\N(h)^3;r,f).
\end{equation}
Put $H_c^3=\min(X,X^2/H^2)$ and $L_b=X/\N(b)^3$. Then
\begin{equation}\label{eq:cubereduction}
 \E(H,X,F;\xi,W)\ll D^\epsilon\left(
 \Sigma\|W\|_{C^{32}}^2+
 \sup_{\N(b)>H_c,\ L_b>1}\E(H,L_b,F;\xi,W)\right).
\end{equation}
The supremum is zero when its set is empty, in particular when $H^2\le X$.
\end{lemma}
\begin{proof}
Expand the right side of \eqref{eq:cubeidentity}; let $c$ be its completion cube index and put $b=hc$. Since $V_*(y)=\sqrt yW(y)$, its norm coefficient becomes $\sqrt{\N(b)}/\sqrt X$. Complete multiplicativity combines all phases into the same whole-index value at $b$, leaving $\sum_{h\mid b}\mu(h)=\ind_{b=1}$. This proves the identity even when $h,c$ or $n,b$ share a prime: a collision zero appears on both sides.

For $\N(h)\le H_c$, apply Cauchy--Schwarz with weight $1/\N(h)$ and Proposition~\ref{prop:completed}. If the fixed polynomial range is $D^{C_0}$, use that proposition with exponent $C_0+1$, since $\N(f)<2F\le2D^{C_0}\le D^{C_0+1}$; its constant changes accordingly. The normalized joint energy is at most
\[
 CD^{\epsilon/3}\|W\|_{C^{32}}^2
 \left(H+H^2FH_c^3/X\right)\le CD^{\epsilon/3}\Sigma\|W\|_{C^{32}}^2,
\]
since $\sum_{\N(h)\le H_c}\N(h)^{-1}\ll\log(2+H_c)$ and $\sum_{\N(h)\le H_c}\N(h)^2\ll H_c^3$. If $H_c<1$ the short sum is empty. All short completion scales are at least $1$.

For the long part, regroup by $b=hc$ to obtain exactly
\[
 \sum_{\N(b)>H_c}\frac{\beta_0(b,f)\chi_b(r)^3}{\N(b)\sqrt{L_b}}
 \sfstar_n a_\xi(n)\chi_n(r)\chi_n(f)^4W(\N(n)/L_b),
\]
where
\[
 \beta_0(b,f)=\bar\alpha(b)^3\xi(b)^3\ind_{(b,f)=1}
 \sum_{h\mid b,\ \N(h)>H_c}\mu(h),\qquad |\beta_0|\le\tauD(b).
\]
Support implies $\N(b)^3\le vX$. Put
\[
 B_X=\sum_{\N(b)^3\le vX}\frac{\tauD(b)}{\N(b)}
 \ll_I\log^2(2+X).
\]
If $P_{\mathrm{long}}(r,f)$ is the normalized long column sum, weighted Cauchy--Schwarz gives
\[
 \frac1F\sfstar_{F\le\N(f)<2F}\sum_{0<\N(r)\le H}|P_{\mathrm{long}}(r,f)|^2
 \le B_X\sum_{\substack{\N(b)>H_c\\\N(b)^3\le vX}}
 \frac{\tauD(b)}{\N(b)}\E(H,L_b,F;\xi,W).
\]
Taking a supremum over the child energies costs at most $B_X^2\ll_I\log^4(2+X)$, absorbed in $D^\epsilon$. If $L_b\le1$ and the term is nonempty, then $L_b\ge1/v$, and ideal counting gives $\E(H,L_b,F)\ll_IH\|W\|_\infty^2\ll_I\Sigma\|W\|_\infty^2$. The remaining terms are precisely the supremum in \eqref{eq:cubereduction}. When $H^2\le X$, $H_c=X^{1/3}$, so a long term cannot have $L_b>1$. No $L$-function of $\xi^3$ is used; $\xi(b)^3$ is only a bounded exterior coefficient.
\end{proof}

\begin{proposition}[Canonical energy with a positive gap]\label{prop:canonical}
Fix $\kappa>0$, $C_0\ge1$, $\epsilon>0$, and let $n_0=\lceil C_0/\kappa\rceil$. Put $J_0=1$ and $J_{j+1}=\max(32,4J_j+12)$. For $\Sigma=XF\le D^{C_0}$, $H,X,F\ge1$ and $H\le\Sigma D^{-\kappa}$,
\begin{equation}\label{eq:canonical}
 \E(H,X,F;\xi,W)\ll_{k,T,S,I,\kappa,C_0,\epsilon}
 D^\epsilon\Sigma\|W\|_{C^{J_{n_0}}(I)}^2.
\end{equation}
In particular, $J_j=36\cdot4^{j-1}-4$ for $j\ge1$.
\end{proposition}
\begin{proof}
Induct on the additional range $H\le D^{j\kappa}$. At $j=0$, ideal counting and the triangle inequality give $\E\ll_IHX\|W\|_\infty^2\ll_I\Sigma\|W\|_\infty^2$. At the next stage use Lemma~\ref{lem:cube} with loss $\epsilon/3$. If its supremum is nonempty, $H^2>X$ and every long length satisfies $L_b<H^2/X$. Apply Proposition~\ref{prop:transfer} at that length with loss $\epsilon/3$. Its children satisfy
\[
 H'\le\frac{HL_b}{\Sigma F}<H(H/\Sigma)^2\le HD^{-2\kappa}\le D^{j\kappa},\quad
 H'/\Sigma'\le D^{-\kappa},\quad \Sigma'\le L_b\le\Sigma.
\]
All child scales are at least $1$ by \eqref{eq:childpositive}. Apply the inductive hypothesis with loss $\epsilon/3$. The three losses add to $\epsilon$. The derivative order required is $\max(32,4J_j+12)$, and the support interval enlarges to $[u/16,4v]$ a fixed finite number of times. Take maxima over these intervals and over the finite twists, after $k,T,S,\kappa,C_0,\epsilon,I$ have been fixed. Finally $H\le D^{C_0}\le D^{n_0\kappa}$, so the finite induction covers the stated range.
\end{proof}

\section{Initialization and the zero-free half-plane}\label{sec:mainproof}
\begin{proposition}[M\"obius--sextic mean square, positive angular type]\label{prop:mobius}
Fix the embedding of $K$ above and a character $\eta=\nu\alpha^k$ with $k\ge1$. For $0<\vartheta\le1/10$, $\epsilon>0$ and $I\Subset(0,\infty)$, some finite $J$ satisfies, for all $D\ge2$ and $W\in C_c^\infty(I)$,
\begin{equation}\label{eq:mobiusbound}
 \sum_{0<\N(u)\le D^{1+\vartheta}}
 \left|\sum_n\mu(n)\eta(n)\chi_n(u)W(\N(n)/D)\right|^2
 \ll_{\eta,S,I,\vartheta,\epsilon}D^{2+\vartheta+\epsilon}\|W\|_{C^J(I)}^2.
\end{equation}
This mean-square statement is proved for positive type only. Negative types in Theorem~\ref{thm:A} follow by analytic conjugation.
\end{proposition}
\begin{proof}
Put $H=D^{1+\vartheta}$ and majorize by $\Phi(\N(u)/H)$. Conjugate the original inner sum before expanding it, and write $n_i=gz_i$. The common phase is now $\tau(z_1)\bar\tau(z_2)$ with $\tau=\alpha^{-k}$. Apply row Poisson to the finite character $\bar\chi_{z_1}\chi_{z_2}$ with exclusion $g$ and use \eqref{eq:pair2}, expanding only $\bar\nu(z)G(z^{-1})$ on $T$. This enters exactly $\mathscr C_k$. Every subsequent conversion stays in that coset. The theta mode remains $m=k+1\ge2$ and never passes through mode zero.

Normalize the smoothed mean by $D^{-1}$ and denote it by $\mathcal M_D$. Its zero frequency is $O_I(H\|W\|_\infty^2)$. Let $e\mid g$ come from the row exclusion and $v\mid(z_1,z_2)$ from coprimality inversion; put $z_i=vm_i$. Common factors $g,v$ cancel by unit modulus before any regrouping. Use the bijection
\[
 b=g/e,\quad f=ev,\quad r=eh;\qquad
 e=(f,r),\quad v=f/e,\quad g=be,\quad h=r/e.
\]
The inverse uses $(v,h)=1$. Both $b,f$ are squarefree and $(b,f)=1$. Since $he^5v^4=rf^4$, the only residual column exclusion is $(m_i,b)=1$; exclusions at $e,v$ are already supplied by the sextic factors. For each finite Fourier term the nonzero contribution is exactly
\begin{equation}\label{eq:initform}
 \frac H{D^2}\sum_{(b,f)=1}^{*}\mu(f)\N(b)\sum_{r\ne0}
 \sum_{(m_1m_2,b)=1}^{*}a_{\tau\chi}(m_1)\overline{a_{\tau\chi}(m_2)}
 \chi_{m_1}(rf^4)\overline{\chi_{m_2}(rf^4)}\mathcal K_{b,f,r},
\end{equation}
where
\[
 \mathcal K_{b,f,r}=W_0(\N(bfm_1)/D)\overline{W_0(\N(bfm_2)/D)}
 \wt\Phi\left(\frac{H\N(r)}{\N(f)^2\N(m_1)\N(m_2)}\right),\quad
 W_0(x)=x^{-1/2}\bar W(x).
\]
These formulas display the angular orientation, weight and all collision masks explicitly.

On $\N(b)\asymp B$, $\N(f)\asymp F$, set
\[
 X=D/(BF),\quad \mathcal H=C_{I,\Phi}D^2/(HB^2),\quad
 \Sigma=XF=D/B,\quad C_{I,\Phi}=\max(2,C_\Phi b_0^2),
\]
where $I=[a_0,b_0]$. Kernel support restricts the rows to $\N(r)\le\mathcal H$ and
\[
 \mathcal H/\Sigma\le C_{I,\Phi}D^{-\vartheta}/B\le C_{I,\Phi}D^{-\vartheta}.
\]
For sufficiently large $D$, depending only on $I,\Phi,\vartheta$, apply Proposition~\ref{prop:canonical} with $\kappa=\vartheta/2$, $C_0=2$. Exclusion removal preserves $\Sigma$ and the gap. If a common-test column scale is below $1$ but its enlarged support is nonempty, it is bounded below by $1/(2b_0)$ and ideal counting supplies the same envelope; if $\mathcal H<1$, the nonzero row sum is empty. Thus all scales are covered.

Set $J_*=J_{n_0}$, $n_0=\lceil4/\vartheta\rceil$, with the sequence of Proposition~\ref{prop:canonical}, and put $q_*=2J_*+4$. For each fixed $b$ we have, for \emph{every} common test $U$ in $I_*=[a_0/8,2b_0]$, a squared-sum bound
\[
 \sum_{f,r}\left|\sfstar_{(n,b)=1}a_{\tau\chi}(n)\chi_n(rf^4)U(\N(n)/X)\right|^2
 \ll D^{\epsilon/2}(XF)^2\|U\|_{C^{J_*}(I_*)}^2,
\]
The rescaled coupled kernels have support in $[a_0/4,b_0]^2$ and satisfy
\[
 \sup_{b,f,r}\|\mathcal K_{b,f,r}\|_{C^{q_*}}
 \le K_{I,\Phi,J_*}\|W\|_{C^{q_*}(I)}^2.
\]
Indeed, their Fourier argument is bounded on nonzero support, their dyadic rescale factors lie in $[1,4)$, and the column variables stay in a fixed positive box. The chain rule expresses each derivative as a finite sum of derivatives of $W_0$, Euler derivatives of $\widehat\Phi$ and bounded rational factors. The constant $K_{I,\Phi,J_*}$ may depend on the fixed derivative order and is independent of $b,f,r,D$. Apply Lemma~\ref{lem:rowkernel} and take $J=q_*$. The contribution of the block to \eqref{eq:initform} is bounded by the envelope times
\[
 \frac{HB}{D^2}\,B\,(XF)^2=H.
\]
The factor $B$ counts possible $b$. Absorb the finite twists and the $O_I(\log^2D)$ blocks into the remaining small power. Then $\mathcal M_D\ll D^\epsilon H\|W\|_{C^J}^2$. Multiplying back by $D$ proves \eqref{eq:mobiusbound}; bounded $D$ is absorbed by counting.
\end{proof}

\begin{lemma}[Sixth-power extraction]\label{lem:extraction}
For every $\delta>0$ there is a finite $J_\delta$ such that
\begin{equation}\label{eq:mobiusindividual}
 \left|\sum_{(n,S)=1}\mu(n)\eta(n)W(\N(n)/D)\right|
 \ll_{\eta,S,I,\delta}\|W\|_{C^{J_\delta}(I)}D^{11/12+\delta}.
\end{equation}
\end{lemma}
\begin{proof}
Write $A_r(D)$ for the inner sum in \eqref{eq:mobiusbound}, and put $Y=D^{(1+\vartheta)/6}$. At good primary primes $p$ with $Y/2<\N(p)\le Y$,
\[
 \chi_n(p^6)=\ind_{p\nmid n},\qquad
 |A_{p^6}(D)-A_1(D)|\ll_I\|W\|_\infty D/Y.
\]
Indeed the difference has only $p$-divisible ideals, whose count is $O_I(D/\N(p))$; if the implied upper norm bound is below $1$, there are none. The prime ideal theorem supplies at least $Y/(4\log Y)$ distinct primes once $Y$ is large. Their primary sixth powers are distinct rows of norm at most $D^{1+\vartheta}$: equality of two sixth powers makes their ratio a sixth root of unity, and the primary congruence forces that unit to be $1$. Averaging $|A_1|^2\le2|A_{p^6}|^2+2|A_{p^6}-A_1|^2$ gives
\[
 |A_1(D)|^2\ll\|W\|_{C^J}^2
 \left(D^{2+\vartheta+\epsilon}Y^{-1}\log Y+D^2Y^{-2}\right).
\]
The square-root exponents are $11/12+5\vartheta/12+\epsilon/2$ and $5/6-\vartheta/6$. Choose $\vartheta=\min(1/10,3\delta/5)$, $\epsilon=\delta/2$, and absorb $\sqrt{\log Y}$ into $D^{\delta/4}$. This yields \eqref{eq:mobiusindividual}, after increasing the constant for bounded $D$. All character parameters and losses were fixed before $D$ varies.
\end{proof}

\begin{proof}[Proof of Theorem~\ref{thm:A}]
By Lemma~\ref{lem:factor} it suffices to consider $\eta=\nu\alpha^k$, $k>0$. Hecke theory makes $L_K(s,\eta)$ entire. Let $L_K^S=L_K\prod_{p\in S}(1-\eta(p)\N(p)^{-s})$, with ramified values zero; its extra factors have no zero for $\Rea s>0$.

Suppose $L_K(\rho,\eta)=0$ with $\Rea\rho>11/12$. Choose $0\ne\varphi\in C_c^\infty((1,2))$, $\varphi\ge0$, and fix $W_\rho(y)=y^{-\rho}\varphi(y)$. Then $\wt W_\rho(\rho)=\int\varphi(y)\dd y/y>0$. For the associated $A_1(D)$, define
\[
 M(s)=\int_0^\infty A_1(D)D^{-s}\dd D/D.
\]
There are no terms for $D<1/2$. Lemma~\ref{lem:extraction} shows this integral is holomorphic on $\Rea s>11/12$: on each compact subset choose $\delta$ smaller than its distance from the boundary, obtaining an integrable majorant, also after any number of $s$ derivatives. In $\Rea s>1$, absolute convergence gives
\[
 M(s)=\wt W_\rho(s)\sum_{(n,S)=1}\mu(n)\eta(n)\N(n)^{-s}
 =\wt W_\rho(s)/L_K^S(s,\eta).
\]
The identity theorem extends $L_K^S(s,\eta)M(s)=\wt W_\rho(s)$ to the entire claimed half-plane. At $s=\rho$ this contradicts the strictly positive right side. This proves holomorphy of the reciprocal on the stated half-plane. Negative types follow by conjugation, and imaginary and real norm shifts by Lemma~\ref{lem:factor}.
\end{proof}

\begin{remark}[The remaining infinity types]\label{rem:radial}
A character with $k=0$ is a finite ray character times $\N^{u+it_0}$. The finite-order theorem of the September 30, 2026 companion paper \cite[Theorem 1.1]{QRH78} supplies the entire $k=0$ case and gives the stronger boundary $u+7/8$, with the principal pole at $1+u+it_0$ permitted. Together with Theorem~\ref{thm:A}, this covers every continuous Hecke character of $K$ of nontrivial infinity type, with its norm shift made explicit. For $a\ne b$, its unitary angular part is algebraic after multiplication by $\N^{|a-b|/2}$.
\end{remark}

\begin{corollary}[CM newforms]\label{cor:newforms}
Let $f$ be a newform of weight $w\ge2$ with complex multiplication by $K=\Q(\sqrt{-3})$, arising by Hecke induction from a primitive character $\psi$ of $K$ of type $(w-1,0)$ or $(0,w-1)$. Then its arithmetic $L$-function has no zeros in
\[
 \Rea s>(w-1)/2+11/12.
\]
The same assertion, with the indicated norm shift, holds for every Hecke character of $K$ covered by Theorem~\ref{thm:A}.
\end{corollary}
\begin{proof}
The arithmetic induction identity is $L(s,f)=L_K(s,\psi)$, with level $N_f=3\N\mathfrak f_\psi$; see \cite[Section 12.3, Theorem 12.5 and its following remark, pp.~211--214]{Iwaniec}. That theorem uses the unitary character $\psi\N^{-(w-1)/2}$ and its Fourier coefficients weighted by $\N(\mathfrak a)^{(w-1)/2}$; primitivity of $\psi$ gives the newform of level $3\N\mathfrak f_\psi$ and all arithmetic local factors. Define $L_{\mathrm{unit}}(s,f)=L(s+(w-1)/2,f)$. Thus
\[
 L_{\mathrm{unit}}(s,f)=L_K(s,\psi\N^{-(w-1)/2}).
\]
Apply Theorem~\ref{thm:A} and restore the norm shift.
\end{proof}

\section{The CM identification for every sextic twist}\label{sec:CM}
We state Deuring's theorem in the all-place form that is needed. If $E/\Q$ has complex multiplication by an order in an imaginary quadratic field $K$, there is a primitive algebraic Hecke character $\psi_E$ of $K$ of type $(1,0)$ (or its conjugate after changing embedding) such that
\begin{equation}\label{eq:Deuring}
 L(E/\Q,s)=L_K(s,\psi_E).
\end{equation}
At a finite place, the Hecke factor is $(1-\psi_E(\mathfrak p)\N\mathfrak p^{-s})^{-1}$ if the local character is unramified, and $1$ otherwise. This statement includes ramified and bad places. The precise reference is \cite[Theorem 4.5, and the local definition following Theorem 4.4, pp.~26--27]{Zhang}; see also \cite[Chapter II, Theorem 10.5]{Silverman} and \cite[Theorem 2.2]{PT}. Panidapu--Thorner use unitary normalization: their $\xi$ is our $\eta_D$, and their $L(s,E)$ is our $L(E_D,s+1/2)$. The product of two Hecke functions in Zhang's Theorem 4.4 concerns $E/K$; over $\Q$ the identity is the single Hecke function \eqref{eq:Deuring}.

\begin{lemma}[Explicit character and normalization]\label{lem:CM}
For $D\in\Q^\times$, the unitary CM character $\eta_D=\psi_D\N^{-1/2}$ has angular type $+1$. On primary ideals prime to $6D$ and its denominators,
\begin{equation}\label{eq:etaD}
 \eta_D(n)=\overline{\chi_n(4D)}\alpha(n).
\end{equation}
The explicit formula agrees with \cite[Chapter II, Example 10.6, pp.~177--178]{Silverman}, whose generator $\pi\equiv2\pmod3$ is $-n$ in our convention. Its primitive conductor is finite, and
\begin{equation}\label{eq:ellipticshift}
 L(E_D,s)=L_K(s,\psi_D)=L_K(s-1/2,\eta_D)
\end{equation}
with all Euler factors included.
\end{lemma}
\begin{proof}
A rational sixth-power scaling replaces $D$ by a sixth-power-free nonzero integer $d$, giving a $\Q$-isomorphism. The endomorphism $(x,y)\mapsto(\omega x,y)$ gives $\operatorname{End}(E_{d,\bar\Q})=\OO_K$, and all its endomorphisms are defined over $K$. Deuring therefore applies with type $(1,0)$ in the chosen embedding. The finite twist is explicit: take $u^6=d$ and $\kappa_d(\sigma)=\sigma(u)/u\in\mu_6$ for $\sigma\in G_K=\operatorname{Gal}(\overline\Q/K)$. The isomorphism
\[
 E_d\longrightarrow E_1,\quad(x,y)\longmapsto(x/u^2,y/u^3)
\]
has conjugate obtained by the CM automorphism $(x,y)\mapsto(\kappa_d^{-2}x,\kappa_d^{-3}y)$. On the type $(1,0)$ eigenspace the twist is $\psi_d=\psi_1\kappa_d^{-1}$; the inverse direction is fixed explicitly by the Jacobi-sum computation below. We use arithmetic Frobenius in Artin reciprocity, so $\kappa_d((\pi))=(d/\pi)_6$ at unramified primes. The Kummer character has finite conductor supported over $6d$, so the infinity type is unchanged. The conductor of the product means its primitive conductor; cancellation at $2$ or $3$ is allowed.

For completeness we derive the good-prime formula, including its direction. At a split rational prime $p\nmid6d$, choose its primary generator $\pi$ and the sextic character $\chi=(\cdot/\pi)_6$ of $\mathbb F_p$. Put $\rho=\chi^3$, $\zeta=\chi^2$ and $J(\zeta,\rho)=\sum_x\zeta(x)\rho(1-x)$. Input~\ref{inp:gauss} and the Gauss--Jacobi relation imply
\[
 J(\zeta,\rho)=\sqrt p\,\gamma_2\gamma_3/\gamma_5
 =-\pi\overline{\chi(4)}.
\]
Indeed $\gamma_5=\chi(-1)/\gamma_1$, $\gamma_1\gamma_2=-\alpha(\pi)\overline{\chi(4)}\gamma_3$, and $\gamma_3^2=\chi^3(-1)$; the quotient $\chi^3(-1)/\chi(-1)=\chi^2(-1)=1$. The conjugate Jacobi sum is $-\bar\pi\chi(4)$. The number of cube roots of $t$ is $1+\zeta(t)+\bar\zeta(t)$, including $t=0$ with the principal term interpreted as $1$. Therefore
\begin{align*}
 \sum_x\rho(x^3+d)
 &=\rho(d)\bigl(\zeta(-d)J(\zeta,\rho)+\bar\zeta(-d)J(\bar\zeta,\rho)\bigr)\\
 &=-\overline{\chi(4d)}\pi-\chi(4d)\bar\pi.
\end{align*}
The missing unweighted character sum is $\sum_t\rho(t+d)=0$. Thus the trace of Frobenius is $\overline{\chi_\pi(4d)}\pi+\chi_\pi(4d)\bar\pi$. The type $(1,0)$ Frobenius belongs to the prime $\pi$: its differential is zero modulo $\pi$, so it is the displayed multiple of $\pi$, rather than its conjugate. This identifies $\psi_d(\pi)$ itself.

For an inert $p\nmid6d$, the cube map on $\mathbb F_p$ is bijective, so $a_p=0$. Its square Frobenius is $-p$ by the characteristic polynomial $T^2+p$. The primary generator of $(p)$ is $-p$, of norm $p^2$; every rational nonzero $x$ satisfies $(x/(-p))_6=1$, since $(p^2-1)/6$ is a multiple of $p-1$. Hence the same character formula gives $\psi_d((p))=-p$ and the correct inert factor $(1+p^{1-2s})^{-1}$. Multiplicativity proves \eqref{eq:etaD} on all good primary ideals. A rational sixth-power change does not alter its sextic symbol outside the enlarged finite bad set.

One can see the seed finite normalization directly. For primary $n$ prime to $2$, cubic reciprocity with the primary element $-2$ gives $\chi_n(4)=(2/n)_3=(n/(2))_3$. Thus $\overline{\chi_n(4)}n$ is the unique generator congruent to $1\pmod{(2)\lambda}$: its $2$-residue is $1$ and its $\lambda$-residue is $1$. The six units give all units modulo $(2)\lambda$. This agrees with the seed conductor $\mathfrak f_1=(2)\lambda$, $\N(E_1)=36$.

At $2,3$ and primes dividing $d$, use the all-place Deuring character, equivalently the primitive character inducing the above finite twist. The induction conductor identity is
\[
 N(E_d)=3\N(\mathfrak f_d),\qquad
 v_{(2)}(\mathfrak f_d)=c_2(E_d)/2,\qquad
 v_\lambda(\mathfrak f_d)=c_3(E_d)-1,
\]
where $c_p(E_d)$ denotes the exponent of $p$ in the elliptic conductor. These identities include primitive conductor cancellation at the wild places. For example $4\cdot16=2^6$ makes the good primary formula for $\psi_{16}$ simply $n\mapsto n$. This character extends to all ideals prime to $3$ and has conductor $\lambda^2$: the modulus $\lambda$ cannot suffice, since $\omega\equiv1\pmod\lambda$ but its type value is $\omega\ne1$. Hence $E_{16}$ has conductor $27$. The model $y^2+y=x^3$, obtained from $E_{16}$ by $x=4X$, $y=8Y+4$, has discriminant $-27$ and is good at $2$, with factor $(1+2^{1-2s})^{-1}$. Deuring's local definition proves \eqref{eq:ellipticshift} also at these places. Finally division by $\N^{1/2}$ makes the character unitary and changes the variable by $1/2$; $\GC((s-1/2)+1/2)=\GC(s)$ checks the gamma normalization. The unitary reflection $s\mapsto1-s$ becomes the arithmetic reflection $s\mapsto2-s$.
\end{proof}

For every $D=dr^6\in\Q^\times$ as in Lemma~\ref{lem:CM}, write
\[
 \mathfrak f_D=\operatorname{cond}(\psi_D)
 =\operatorname{cond}(\eta_D)=\mathfrak f_d,\qquad
 N_D=N(E_D)=3\N\mathfrak f_D.
\]

\begin{proof}[Proof of Theorem~\ref{thm:B}]
Apply Theorem~\ref{thm:A} to \eqref{eq:ellipticshift}: $1/2+11/12=17/12$. Alternatively, the good-prime identity suffices for this zero-free half-plane. Removing the finite bad set and putting its factors back changes the elliptic function only by rational local factors whose Frobenius roots have modulus at most $p^{1/2}$, and changes the unitary Hecke function by factors with roots of modulus at most $1$. They have no zero or pole in the claimed region. Thus finitely many bad Euler factors cannot affect the conclusion.

The completed function
\[
 \Lambda_D(s)=N_D^{s/2}(2\pi)^{-s}\Gamma(s)L(E_D,s)
\]
is entire and satisfies $\Lambda_D(s)=w_D\Lambda_D(2-s)$, $w_D\in\{\pm1\}$, by \cite[Chapter II, Corollary 10.5.1]{Silverman}. Nontrivial zeros are the zeros of $\Lambda_D$, equivalently the zeros of $L(E_D,s)$ other than $s=0,-1,-2,\ldots$, where the poles of $\Gamma(s)$ are cancelled. The Euler product is absolutely convergent and nonzero for $\Rea s>3/2$, so a nontrivial zero has real part at most $3/2$ and, by reflection, at least $1/2$. The zero-free half-plane just proved and the same reflection therefore confine these zeros to $7/12\le\Rea s\le17/12$. The boundary lines remain permitted because the zero-free half-planes are open.
\end{proof}

\section{Power-saving applications}\label{sec:applications}
We establish frequency uniformity anew from the analytic conductor. We will also use the finite-order $7/8$ theorem of \cite[Theorem 1.1]{QRH78}, with the principal pole allowed. These are the only extra release statements used in this section.

\begin{lemma}[Uniform logarithmic derivative and prime sums]\label{lem:uniform}
Let $\eta=\nu\alpha^k$ have primitive conductor $\mathfrak f$, $k\ne0$, and put
\[
 \mathcal C(\eta,t)=3\N\mathfrak f\,(|k|+|t|+3)^2.
\]
For $0<\delta<1/12$ and $11/12+\delta\le\sigma\le2$,
\begin{equation}\label{eq:logderivative}
 |L_K'/L_K(\sigma+it,\eta)|\ll_\delta\log\mathcal C(\eta,t).
\end{equation}
The constant is independent of $k$ and $\mathfrak f$. For $x\ge2$,
\begin{equation}\label{eq:primeuniform}
 \sum_{\N\mathfrak p\le x}\eta(\mathfrak p)\log\N\mathfrak p
 \ll_\delta x^{11/12+\delta}
 \log\bigl(3\N\mathfrak f\,(|k|+x+3)^2\bigr)\log(2x).
\end{equation}
For the finite-order cases below take $0<\delta<1/8$ and $b=7/8+\delta$. For nonprincipal characters covered by the $7/8$ theorem, the same bounds hold with $11/12$ replaced by $7/8$ and $k=0$, uniformly in the primitive conductor, on $b\le\sigma\le2$. For the primitive principal character the logarithmic-derivative assertion is
\begin{equation}\label{eq:principalderivative}
 \left|-\frac{\zeta_K'}{\zeta_K}(s)-\frac1{s-1}\right|
 \ll_\delta\log\bigl(3(|t|+3)^2\bigr),
 \qquad 7/8+\delta\le\Rea s\le2,
\end{equation}
with removable continuation of the difference at $s=1$; the prime sum is $x$ plus the corresponding $7/8$ error. The degree-one variants use the degree-one analytic conductor and the same pole subtraction; they are used here only for the fixed functions $\zeta(s)$ and $L(s,\chi_{-3})$. Deleted Euler factors add at most $C_\delta\log\N\mathfrak m$ for a presentation modulus $\mathfrak m$.
\end{lemma}
\begin{proof}
Euler convergence on $\sigma=2$ gives an absolute bound for $\log L(2+it)$, with the logarithm fixed by its Euler series. Write $q=3\N\mathfrak f$. On $\sigma=-1$, the functional equation gives
\[
 |L_K(-1+it,\eta)|\ll q^{3/2}(|k|+|t|+3)^3.
\]
This is uniform in $k$: the gamma quotient has modulus $(2\pi)^{-3}|z(z+1)(z+2)|$, where $z=-1+|k|/2+it$, by $\Gamma(z+3)/\Gamma(z)=z(z+1)(z+2)$. The dual Euler series is on $\sigma=2$ and is bounded by $\zeta_K(2)$. On $\sigma=4$, Euler convergence bounds $L$ absolutely. Divide by $(s+|k|+3)^3$ and apply Phragm\'en--Lindel\"of on $[-1,4]$; the denominator has no zero there, and the order-one growth supplied by Hecke theory permits this strip argument. Both boundary bounds are $Cq^{3/2}$, independently of $k$. Consequently
\[
 |L_K(\sigma+it,\eta)|\ll q^{3/2}(|k|+|t|+3)^3,
 \qquad -1\le\sigma\le4,
\]
so $\log|L|\ll\log\mathcal C(\eta,t)$, uniformly. No constant from Proposition~\ref{prop:mobius} enters this estimate.

At center $2+it$ take a disk of radius $2-11/12-\delta/2$. It lies in the zero-free half-plane, so $\log L$ is holomorphic there. Its real part has the bound just obtained, and its center value is bounded by the Euler series. Borel--Carath\'eodory and Cauchy's derivative estimate on a smaller disk containing $11/12+\delta+it$ prove \eqref{eq:logderivative}; points further right are covered by the same disk or the absolutely convergent Euler series. This proof works for arbitrarily large $k$ and conductor because only their logarithmic analytic conductor occurs.

For the principal case apply the strip-growth argument to the entire function $(s-1)L(s)$, dividing by $(s+3)^4$ and increasing the polynomial degree to four. Thus the function $H(s)=(s-1)L(s)/(s+1)$ has polynomial-conductor growth on the positive disks used here; its center logarithm is bounded because $(1+it)/(3+it)$ has modulus between $1/3$ and $1$. Cancel the pole in this way before taking the logarithm. The same argument gives $-L'/L(s)=1/(s-1)+O_\delta(\log\mathcal C)$ in the $7/8+\delta$ strip, with bounded elementary corrections. For a nonprincipal finite-order character use the same gamma quotient with $k=0$: its modulus is $(2\pi)^{-3}|(-1+it)it(1+it)|$, interpreted by continuity when $t=0$. The entire strip-growth argument and the zero-free disks with left boundary $7/8+\delta/2$ therefore give constants independent of the primitive conductor. For primitive Dirichlet characters, the gamma shift is one of the two fixed parities; the functional equation gives a polynomial bound with logarithm $O(\log(q(|t|+3)))$, uniformly in the conductor $q$, and the coefficient bound is $|b(n)|\le\Lambda(n)$. The same disks and principal regularization prove the degree-one assertions.

Here are the details of the sharp prime sum. The Dirichlet coefficients $b(n)$ of $-L_K'/L_K$ satisfy $|b(n)|\le2\Lambda(n)$: at a split prime there are two terms, and at an inert prime one term of norm $p^2$ and weight $2\log p$. Take a half-integer $x\ge e$, with $c=1+1/\log x\le2$ and $T=x$; the range $2\le x<e$ is trivial. Truncated Perron gives
\[
 \sum_{n\le x}b(n)=\frac1{2\pi i}\int_{c-iT}^{c+iT}-\frac{L'}L(s)\frac{x^s}s\dd s+O(\log^2(2x)).
\]
For clarity the error follows from $|b(n)|\le2\log n$ and
$\min(1,(T|\log(x/n)|)^{-1})$: split $n$ into $[x/2,2x]$ and its complement; the near sum is $O(x\log^2x/T)$ because $|x-n|\ge1/2$, and the complementary sum is bounded using Euler convergence at $c$.

Shift to $b=11/12+\delta$ for nonzero angular type and to $b=7/8+\delta$ for the finite-order cases. There are no zeros in the rectangle; its only possible pole is at $1$ in the principal case. The new vertical side is
\[
 O_\delta\bigl(x^b\log\mathcal C(\eta,T)\log(2T)\bigr),
\]
and the horizontal sides are $O_\delta(x^c\log\mathcal C(\eta,T)/T)$. If principal, the sole residue at $1$ is $x$. Prime ideal powers of exponent at least two have total absolute weight $O(\sqrt x\log^2(2x))$, which is absorbed. Moving an arbitrary real $x$ to a half-integer changes the sum by $O(\log(2x))$. This proves \eqref{eq:primeuniform}. If $\mathfrak m$ denotes an imprimitive presentation modulus, its removed Euler factors contribute at most $C_\delta\sum_{p\mid\mathfrak m}\log\N(p)\le C_\delta\log\N\mathfrak m$ to the logarithmic derivative in the positive strip. The primitive bound uses its primitive analytic conductor; a fixed deletion adds only a fixed local contribution.
\end{proof}

\begin{theorem}[Eisenstein prime ideals in sectors]\label{thm:sectors}
Let $\pi_{\rm Eis}(x;J)$ count prime ideals of norm at most $x$, prime to $3$, whose primary generators have arguments in an arc $J\subset\R/2\pi\Z$. For every $\epsilon>0$, uniformly in $J$,
\begin{align}\label{eq:sectorweighted}
 \sum_{\substack{\N\mathfrak p\le x,\ \mathfrak p\nmid3\\\arg\pi_\mathfrak p\in J}}\log\N\mathfrak p
 &=\frac{|J|}{2\pi}x+O_\epsilon(x^{11/12+\epsilon}),\\
 \pi_{\rm Eis}(x;J)&=\frac{|J|}{2\pi}\operatorname{Li}(x)+O_\epsilon(x^{11/12+\epsilon}).\label{eq:sectorcount}
\end{align}
Prime ideals of residue degree two may be omitted with an $O(\sqrt x\log x)$ change in the weighted formula. Counting all six associates of each prime element multiplies the main density accordingly.
\end{theorem}
\begin{proof}
For every nonzero integer $j$, the primary-generator character $\alpha^j$ is a Hecke character of conductor dividing $3\OO$ and angular type $j$. Lemma~\ref{lem:uniform} bounds its weighted prime sum by
$C_\delta x^{11/12+\delta}\log^2(2x)$ uniformly for $|j|\le x^{1/12}$: the conductor is fixed and the gamma parameter is $|j|$. The zero frequency is $x+O_\delta(x^{7/8+\delta}\log^2x)$, by the finite-order theorem for $\zeta_K$.

Apply the Selberg trigonometric majorant and minorant for an arc with degree $M$ \cite[Lemma 2.1 and its proof]{PT}. Their constant coefficients differ from $|J|/(2\pi)$ by $O(M^{-1})$, and their nonconstant coefficients have absolute value at most $C/|j|+C/M$. Rotation, complementation and splitting a wrapping arc extend the interval polynomials to every circle arc. Equivalently this is the weighted Erd\H{o}s--Tur\'an inequality
\[
 \left|\sum w_\mathfrak p\ind_J(\theta_\mathfrak p)-\frac{|J|}{2\pi}\sum w_\mathfrak p\right|
 \ll\frac{\sum w_\mathfrak p}{M+1}
 +\sum_{1\le|j|\le M}\frac1{|j|}\left|\sum w_\mathfrak p e^{ij\theta_\mathfrak p}\right|,
\]
valid for nonnegative weights and any endpoint convention. With $w_\mathfrak p=\log\N\mathfrak p$ and $M=\lfloor x^{1/12}\rfloor$, this is
$O_\delta(x^{11/12}+x^{11/12+\delta}\log^3(2x))$. Choose $\delta<\epsilon/2$ and absorb the logarithms. Partial summation proves \eqref{eq:sectorcount}; the prime ideals of degree two have norms $p^2$ and contribute only $O(\sqrt x\log x)$.
\end{proof}

This sharpens Hecke's equidistribution by a fixed power saving in the full-range discrepancy. Coleman's classical region is \eqref{eq:classicalregion}; combining it with the standard contour argument gives the familiar subexponential full-range errors. Classical results also localize the primes and angles: \cite{ColemanPrimes}, as stated in \cite[Introduction]{Khale}, gives sector asymptotics in imaginary quadratic fields for angular widths larger than $x^{-5/24+\epsilon}$ and interval lengths at least $x^{19/24+\epsilon}$. Those short-interval statements are distinct from the full-range power-saving discrepancy proved above.

\begin{theorem}[The CM Sato--Tate distribution with a power saving]\label{thm:ST}
Fix $D\in\Q^\times$. For every interval $I\subset[-1,1]$, with its membership at $0$ specified, put
\[
 \mu_{\rm CM}(I)=\frac12\ind_{0\in I}+\frac1{2\pi}\int_I\frac{dt}{\sqrt{1-t^2}}.
\]
Then for every $\epsilon>0$, uniformly in $I$,
\begin{align}\label{eq:STweighted}
 \sum_{\substack{p\le x,\ p\text{ good for }E_D\\a_p(E_D)/(2\sqrt p)\in I}}\log p
 &=\mu_{\rm CM}(I)x+O_{D,\epsilon}(x^{11/12+\epsilon}),\\
 \#\{p\le x:p\text{ good},\ a_p(E_D)/(2\sqrt p)\in I\}
 &=\mu_{\rm CM}(I)\operatorname{Li}(x)+O_{D,\epsilon}(x^{11/12+\epsilon}).\label{eq:STcount}
\end{align}
\end{theorem}
\begin{proof}
For every $j\ne0$, $\eta_D^j$ has angular type $j$ and primitive conductor dividing the fixed $\mathfrak f_D$. Hence Lemma~\ref{lem:uniform}, with the same argument as for sectors, gives angular equidistribution of its good prime ideals with weighted discrepancy $O_{D,\epsilon}(x^{11/12+\epsilon})$. Removing the fixed bad set costs logarithmic terms; removing degree-two prime ideals costs $O(\sqrt x\log x)$.

If $p$ splits, its two good prime ideals have conjugate character values $e^{i\theta},e^{-i\theta}$, and $a_p/(2\sqrt p)=\cos\theta$. For $I=[a,b]$, its inverse image under cosine is a union of at most two arcs of total length $2(\arccos a-\arccos b)$. Each rational split prime is counted twice. Its main weighted mass is therefore
\[
 \frac{\arccos a-\arccos b}{2\pi}x
 =\frac x{2\pi}\int_a^b\frac{dt}{\sqrt{1-t^2}}.
\]
Arc majorants also handle open endpoints and degenerate intervals with the same discrepancy. At inert primes the trace is zero. Their weighted count is $x/2+O_\epsilon(x^{7/8+\epsilon})$, by the finite-order $7/8$ theorem for $\zeta(s)$ and $L(s,\chi_{-3})$ and Lemma~\ref{lem:uniform}. A good split prime has nonzero trace: its Frobenius lies in $K$, and a trace-zero algebraic integer there has norm $3r^2$, which cannot be a split prime. Thus the atom is exactly the inert contribution. Adding the two masses proves \eqref{eq:STweighted}. Partial summation proves \eqref{eq:STcount}.
\end{proof}

For the classical quantitative comparison, \cite[Theorem 1.1]{PT} proves that for a fixed CM curve, $\delta,\theta>0$ with $\delta+\theta<5/24$, intervals of primes of length $h\ge x^{1-\delta}$ and angular intervals of length at least $x^{-\theta}$ have relative error
\[
 O\left(\exp\left(-c\frac{(\log x)^{1/3}}{(\log\log x)^{1/3}}\right)\right).
\]
This is a strong unconditional shrinking-sector, short-interval result. The fixed power saving in \eqref{eq:STweighted} concerns the full-range distribution and follows from the common zero-free half-plane.


\begin{lemma}[Conductor and rational height]\label{lem:height}
For $D=a/b\ne0$ in lowest terms, $b>0$, put $H(D)=\max(|a|,b)$ and $N_D=N(E_D)$. There is a signed sixth-power-free integer $d$ with $D=dr^6$, $r\in\Q_{>0}$, and
\[
 |d|\le |a|b^5\le H(D)^6,\qquad
 N_D\le46656|d|^2\le46656H(D)^{12}.
\]
For integer $D$ one also has $|d|\le|D|$ and $N_D\le46656|D|^2$.
\end{lemma}
\begin{proof}
Reduce every $v_p(D)$ modulo six into $\{0,\ldots,5\}$. This gives $d$ and $r$, and the sixth-power scaling identifies $E_D$ with $E_d$ over $\Q$. A numerator prime contributes at most its exponent in $|a|$; a denominator prime of exponent $v$ contributes at most $5v$. These observations prove the bounds for $d$.

By Lemma~\ref{lem:CM}, $\psi_d$ is the seed character, of conductor $(2)\lambda$, times a Kummer character of order dividing six supported over $6d$. At a place of residue characteristic prime to six, this character is trivial on the pro-$p$ group $1+\mathfrak p$, so its conductor exponent is at most one. At either place over six, Hensel's lemma gives
\[
 1+\mathfrak p^{\,2v_{\mathfrak p}(6)+1}
 \subset (K_{\mathfrak p}^{\times})^6:
\]
for a unit $u$ in the left side, apply $v(f(1))>2v(f'(1))$ to $f(X)=X^6-u$. A character of order dividing six is trivial on sixth powers. Thus its conductor exponent is at most three at $(2)$ and five at $\lambda$. Multiplication by the seed preserves these upper bounds, and hence
\[
 \mathfrak f_d\mid(2)^3\lambda^5
 \prod_{\substack{\mathfrak p\mid d\\\mathfrak p\nmid6}}\mathfrak p,
 \qquad \N\mathfrak f_d\le2^6\,3^5|d|^2.
\]
The primes above a rational $p\ge5$ have product of norms $p^2$, both in the split and inert cases. The conductor identity $N_D=3\N\mathfrak f_d$ proves the result. For later use, at a rational prime $p\ge5$ the curve $E_d$ is good exactly when $p\nmid d$: the discriminant $\Delta(E_d)=-432d^2$ of the model $y^2=x^3+d$ has valuation zero in that case, and otherwise its valuation $2v_p(d)\in\{2,4,6,8,10\}$ cannot become zero under a Weierstrass change, which changes discriminant valuations by multiples of twelve.
\end{proof}

\begin{corollary}[The least prime in a CM interval]\label{cor:least}
Fix an interval $I\subset[-1,1]$, with its endpoints specified, and suppose $\mu_{\rm CM}(I)>0$. For every fixed $A>12$, the least prime good for $E_D$ with
\[
 \frac{a_p(E_D)}{2\sqrt p}\in I
\]
satisfies
\begin{equation}\label{eq:leastCM}
 p\ll_{I,A}\bigl(\log(N(E_D)+3)\bigr)^A
 \qquad(D\in\Q^\times).
\end{equation}
The constant is independent of $D$. In particular this bounds the least good prime with $a_p(E_D)<0$. For integer $D$ the right side can be replaced by $C_{I,A}(\log(|D|+2))^A$, and for rational $D$ by $C_{I,A}(\log(H(D)+2))^A$.
If $0\in I$, the stronger bound $p\ll\log(N(E_D)+3)$ holds with an absolute constant.
\end{corollary}
\begin{proof}
Put $L_D=\log(N_D+3)\ge1$, and let $S_D$ consist of the prime ideals over $2,3$ and the bad rational primes. Every prime of bad reduction divides $N_D$, so
\[
 \sum_{\mathfrak p\in S_D}\log\N\mathfrak p
 \le\log12+2\log N_D=O(L_D).
\] The primitive conductors of the powers of $\eta_D$ divide $\mathfrak f_D$. In applying Lemma~\ref{lem:uniform}, delete $S_D$ even when a power becomes unramified there, so that the remaining values are exactly the powers of the original good-prime character.

Suppose $0\notin I$. Positive CM measure then gives a positive-length closed subinterval $I_0\subset I\cap(-1,1)$ separated from zero. The set
\[
 \mathcal J=\{\theta\in\R/2\pi\Z:\cos\theta\in I_0\}
\]
is a union of two arcs of positive normalized Haar measure $v_I=|\mathcal J|/(2\pi)$. Sum the Selberg minorants for these arcs used in Theorem~\ref{thm:sectors}. For a sufficiently large fixed degree $M=M_I$ their sum is a real polynomial
\[
 P(\theta)=c_0+\sum_{1\le|j|\le M}c_j e^{ij\theta}
 \le\ind_{\mathcal J}(\theta),\qquad c_0\ge v_I/2>0.
\]
The two arcs are disjoint, so the sum preserves the pointwise minorant property. Its coefficients and supremum norm depend only on $I$.

Choose
\[
 \delta=\tfrac12(1/12-1/A)>0,\qquad
 b_*=11/12+\delta,\qquad \gamma=1-b_*,
\]
so $A\gamma>1$. For $1\le|j|\le M$, Lemma~\ref{lem:uniform} and conjugation give, uniformly in $D$,
\[
 \left|\sum_{\substack{\N\mathfrak p\le x\\\mathfrak p\notin S_D}}
       \eta_D(\mathfrak p)^j\log\N\mathfrak p\right|
 \ll_{I,A}x^{b_*}(L_D+\log(2x))\log(2x)+L_D.
\]
Indeed $3\N\mathfrak f_D=N_D$ and $j$ is bounded by the fixed $M$. For the zero frequency the principal prime ideal theorem gives $x+o(x)$, with its threshold depending only on $K$; deletion costs $O(L_D)$. Removing degree-two prime ideals costs $O_I(\sqrt x\log(2x))$ even after multiplication by $P$. At a remaining split prime the two character values are conjugate and both cosines are $a_p/(2\sqrt p)$. Pointwise minorization with nonnegative logarithmic weights, followed by division by two, therefore gives
\begin{equation}\label{eq:leastminorant}
 \begin{split}
 W_D(I;x)&:=\sum_{\substack{p\le x,\ p\ {\rm good}\\
                  a_p(E_D)/(2\sqrt p)\in I}}\log p\\
 &\ge c_Ix-C_{I,A}x^{b_*}(L_D+\log(2x))\log(2x)
       -C_{I,A}L_D-C_{I,A}\sqrt x\log(2x),
 \end{split}
\end{equation}
for $x\ge x_I$, where $c_I>0$ and $x_I$ are independent of $D$. The interior interval $I_0$ respects every endpoint convention.

Take $x=C L_D^A$, with $C\ge x_I$. The principal error divided by $x$ is bounded by
\[
 C_{I,A}C^{-\gamma}L_D^{\,1-A\gamma}
       (1+\log C+\log L_D)^2.
\]
Since $A\gamma>1$, the supremum over $L_D\ge1$ tends to zero as $C\to\infty$. The other two errors divided by $x$ also tend uniformly to zero. A sufficiently large $C=C_{I,A}$ makes \eqref{eq:leastminorant} positive and hence supplies a prime.

If $0\in I$, use good inert primes, whose traces are zero. The classical prime number theorem for the fixed progression $p\equiv2\pmod3$, obtained from $\zeta$ and $L(s,\chi_{-3})$ and their classical zero-free regions, gives
\[
 \sum_{\substack{3<p\le x\\p\equiv2(3)}}\log p=x/2+o(x).
\]
Its error is at most $x/4$ for $x\ge x_0$ with absolute $x_0$. Every bad prime divides $N_D$, so deleting the bad rational primes costs at most $\sum_{p\ {\rm bad}}\log p\le\log N_D\le L_D$. Taking $x=C L_D$ with absolute $C\ge\max(x_0,8)$ leaves a positive sum. This also treats the singleton $\{0\}$. Take $I=[-1,0)$ for the negative-trace assertion. Finally Lemma~\ref{lem:height} gives the integer and rational-height formulations.
\end{proof}

For comparison, \cite[Theorem 1.1]{CPS} gives the unconditional Linnik bound $p\ll(N_D/(\beta-\alpha))^{A_0}$ for a prescribed positive-length angle interval $[\alpha,\beta]$, with an absolute effectively computable exponent $A_0>2$. For the family $y^2=x^3+D$ and a fixed interval $I$ of positive CM measure, Corollary~\ref{cor:least} replaces polynomial conductor dependence by polylogarithmic dependence.

\begin{remark}[Distinguishing two sextic twists]\label{rem:distinguish}
The finite-order $7/8$ theorem of \cite[Theorem 1.1]{QRH78}, together with classical CM theory, implies the following separate consequence. If $E_D$ and $E_{D'}$ are not isogenous over $\Q$, then for every fixed $A>8$ some prime good for both satisfies
\[
 a_p(E_D)\ne a_p(E_{D'}),\qquad
 p\ll_A\bigl(\log(N_DN_{D'}+3)\bigr)^A.
\]
In particular one may take exponent $9$. For integer inputs the logarithm may be replaced by $\log(|DD'|+2)$, and for rational inputs by $\log(H(D)H(D')+2)$.

Here is the proof and its dependence. The primitive character
$\xi=\eta_D\eta_{D'}^{-1}=\psi_D\psi_{D'}^{-1}$
has finite order $h\mid6$ and conductor dividing $\mathfrak f_D\mathfrak f_{D'}$.
Deuring's isogeny correspondence \cite[Theorem 4.5]{Zhang} shows that $\xi=1$ would make the curves isogenous over $\Q$; their positive-type characters cannot be interchanged with a conjugate of type $(0,1)$. Thus $2\le h\le6$.
Put $L=\log(N_DN_{D'}+3)$ and remove primes over $2,3$ and bad primes for either curve. Every such bad rational prime divides $N_DN_{D'}$, so the logarithmic cost is at most $\log12+2\log(N_DN_{D'})=O(L)$. On the remaining prime ideals, character orthogonality gives
\[
 \ind_{\xi(\mathfrak p)\ne1}
   =1-\frac1h\sum_{j=0}^{h-1}\xi(\mathfrak p)^j.
\]
Choose $\delta=\tfrac12(1/8-1/A)>0$ and $b_*=7/8+\delta$.
The conductor-uniform finite-order bounds of Lemma~\ref{lem:uniform}, including its principal main term, give
\[
 \sum_{\substack{\N\mathfrak p\le x,\ \mathfrak p\notin S\\
                    \xi(\mathfrak p)\ne1}}\log\N\mathfrak p
 =\left(1-\frac1h\right)x
  +O_A\bigl(x^{b_*}(L+\log(2x))\log(2x)+L\bigr).
\]
The constant is independent of both curves since $h\le6$.
Removing degree-two prime ideals costs $O(\sqrt x\log(2x))$.
As $A(1-b_*)>1$, the same estimate used in \eqref{eq:leastminorant}, with $x=C_A L^A$, makes the remaining sum positive. It supplies a split prime good for both with $\xi(\mathfrak p)\ne1$.

For $\mathfrak p=(\pi)$ above this $p\ge5$, write
$z=\psi_D(\mathfrak p)=u\pi$ and
$z'=\psi_{D'}(\mathfrak p)=u'\pi$, $u,u'\in\mu_6$.
These formulas hold on good primes by the primitive CM character: choose sixth-power-free representatives of $D,D'$ so $p$ is prime to their parameters. Both numbers generate $\mathfrak p$ and have norm $p$.
If the traces were equal, $z,z'$ would be roots of the same polynomial $T^2-a_pT+p$, so $z'=z$ or $z'=\bar z$. The first contradicts $\xi(\mathfrak p)\ne1$; the second contradicts $(z')=\mathfrak p\ne\overline{\mathfrak p}=(\bar z)$.
Thus the traces differ. Lemma~\ref{lem:height} gives the height and integer formulations.
This proof uses the finite-order companion theorem for every nonzero frequency of $\xi$.
\end{remark}

\begin{theorem}[Every symmetric power]\label{thm:sym}
For every integer $n\ge1$, the arithmetic symmetric-power function of $E_D$ is zero-free in
\[
 \Rea s>n/2+11/12.
\]
If $n=2r$ is even, its finite-order middle factor is zero-free in $\Rea s>r+7/8$. The possible principal pole at $s=r+1$ is permitted. More precisely, put $L_n^{\mathrm{unit}}(s)=L(\operatorname{Sym}^nE_D,s+n/2)$ and $\chi_K=\chi_{-3}$. Then
\begin{align}\label{eq:symodd}
 L_{2r+1}^{\mathrm{unit}}(s)
 &=\prod_{j=0}^{r}L_K(s,\eta_D^{2r+1-2j}),\\
 L_{2r}^{\mathrm{unit}}(s)
 &=L_\Q(s,\chi_K^r)\prod_{j=0}^{r-1}L_K(s,\eta_D^{2r-2j}).\label{eq:symeven}
\end{align}
In these formulas $\eta_D^m$ denotes the primitive character inducing the power, including any cancellation of finite ramification. The power $\chi_K^r$ is also taken in the primitive convention, so it is the trivial character of conductor one when $r$ is even. The case $n=0$ is $\zeta(s)$ and has the $7/8$ boundary of the finite-order theorem.
\end{theorem}
\begin{proof}
Put $V_\ell=T_\ell(E_D)\otimes\Q_\ell$ and diagonalize the $\OO_K$-action on $V_\ell\otimes\overline{\Q}_\ell$. The subgroup $G_K$ fixes its two eigenspaces and the other coset exchanges them, giving induction from $\psi_D$. The nondegenerate Galois-equivariant Weil pairing identifies $\bigwedge^2V_\ell$ with $\Q_\ell(1)$, so its determinant is the cyclotomic character $\operatorname{cyc}$. Use a basis $e_1,e_2$ on which restriction to $G_K$ has the two conjugate characters. In $\operatorname{Sym}^n$, pair the monomials $e_1^{n-j}e_2^j$ and $e_1^je_2^{n-j}$ for $j<n/2$. Their restriction has characters $\psi_D^{n-2j}\operatorname{cyc}^j$ and its conjugate, since $\psi_D\bar\psi_D=\operatorname{cyc}$ on $G_K$. Each pair is the induction of the former character. For even $n=2r$, the remaining line $e_1^re_2^r$ restricts to $\operatorname{cyc}^r$. An element outside $G_K$ interchanges the two basis vectors. If its off-diagonal entries are $a,b$, then $ab=-\det V_\ell=-\operatorname{cyc}$, so its value on this line is $(-1)^r\operatorname{cyc}^r$. The line is therefore $\operatorname{cyc}^r\chi_K^r$.

This is a decomposition of representations, so it respects inertia invariants and all finite local factors, including bad primes. After the norm shifts each paired factor is $L_K(s,\eta_D^{n-2j})$, proving \eqref{eq:symodd}--\eqref{eq:symeven}. All paired factors have positive nonzero angular type and are zero-free for $\Rea s>11/12$ by Theorem~\ref{thm:A}. The remaining character is either $\chi_K$ or the principal Dirichlet character, so \cite[Theorem 1.1]{QRH78} gives $7/8$ for it. A product of these nonzero meromorphic factors has no zero in the stated region. Restoring the arithmetic shift gives $n/2+11/12$.
\end{proof}

Classically the same CM decomposition gives analytic continuation and shrinking zero-free regions for each fixed symmetric power from Hecke theory and Coleman's region, together with classical regions for the finite-order factors. The $7/8$ input here is restricted to the latter factors; no $7/8$ assertion for a nonzero angular type is used.

\section{The limits of the hypotheses}\label{sec:limits}
The scalar theta mode $m=0$ retains the height constant $Cv^{2/3}$. Its inverse-cusp term $C_0v^{-2/3}$ contributes $C_0\int_0^1v^{2s-8/3}\,dv=C_0/(2s-5/3)$, giving a pole at $s=5/6$ when $C_0\ne0$; the reflected Mellin function has the partner at $s=1/6$. This is the scalar polar term in \cite[Proposition 5.1]{DR}. Residue-free reflection uses the vanishing of these constants in the pure jets. The positive-type initialization keeps $m=k+1\ge2$, and conjugation handles negative types without entering that scalar mode. Original angular type $k=0$ instead gives p2's mode $m=1$: the principal target has a pole at $s=1$, where its reciprocal has a zero. That is compatible with zero-freeness of the $L$-function and must not be read as absence of poles.

Complete multiplicativity and unit modulus suffice for transfer, but the completed estimate additionally needs a finite Hecke mask and an integer angular theta mode. For the ideal Liouville function $\lambda_K(\mathfrak a)=(-1)^{\Omega(\mathfrak a)}$,
\[
 L(s,\lambda_K)=\zeta_K(2s)/\zeta_K(s)
\]
has a zero at $1$, and $\mu\lambda_K=\mu^2$ has a positive density main term. The finite-conductor periodic-mask hypothesis in Proposition~\ref{prop:completed} excludes this function. For a unitary Hecke character $\eta$ of nonzero angular type let $\eta^{[2]}$ be the primitive character inducing its square, and put
\[
 \mathcal P_\eta=\{\mathfrak p:\mathfrak p\mid\mathfrak f_\eta,\ 
                      \mathfrak p\nmid\mathfrak f_{\eta^{[2]}}\}.
\]
With $\eta$ extended by zero at primes dividing $\mathfrak f_\eta$,
\begin{equation}\label{eq:Liouville}
 L(s,\lambda_K\eta)=
 \frac{L_K(2s,\eta^{[2]})}{L_K(s,\eta)}
 \prod_{\mathfrak p\in\mathcal P_\eta}
 (1-\eta^{[2]}(\mathfrak p)\N\mathfrak p^{-2s}).
\end{equation}
At an originally unramified prime the local identity is
$(1-z)/(1-z^2)=(1+z)^{-1}$, with
$z=\eta(\mathfrak p)\N\mathfrak p^{-s}$.
At a prime in $\mathcal P_\eta$ the displayed factor cancels the newly restored primitive-square Euler factor; at primes where both characters are ramified all factors are $1$.
Thus \eqref{eq:Liouville} follows for $\Rea s>1$ by Euler products and everywhere by meromorphic continuation.
For $\Rea s>1/2$ its numerator is nonzero by absolute Euler convergence, every finite factor is nonzero, and the denominator is entire. The meromorphic twist has no zeros there and can have poles at denominator zeros.
For $\lambda_K$ the zero at $s=1$ comes from the pole of $\zeta_K(s)$; for nonzero angular type the denominator $L_K(s,\eta)$ is entire, which is why the preceding argument applies.

An integer circle frequency supplies both a multiplicative angular character and a pure derivative of integral order. The Euler product identifies the M\"obius Mellin transform with the reciprocal of the Hecke function.

Finally a real norm power is a translation of the variable, not a unitary phase: $L(s,\eta\N^u)=L(s-u,\eta)$. Both the coefficient bounds and the zero-free boundary must be shifted. These restrictions explain exactly why the argument does not extend its unshifted conclusion to arbitrary multiplicative functions or nonunitary characters.

\appendix
\section{Smooth kernels and recombination}\label{app:smooth}
We record the statements used from p2 Appendix B with their coefficient-independent proofs.
\begin{lemma}[Mellin separation; p2 Lemma B.1]\label{lem:smooth}
For a fixed box $\mathcal I=I_1\times\cdots\times I_d\Subset(0,\infty)^d$, a smooth kernel supported there has
\[
 \mathcal K(\boldsymbol x)=\int_{\R^d}b(\boldsymbol t)\prod_jx_j^{it_j}\dd\boldsymbol t,
 \qquad
 \int|b(\boldsymbol t)|(1+|\boldsymbol t|)^J\dd\boldsymbol t
 \ll\|\mathcal K\|_{C^q(\mathcal I)}
\]
for every even integer $q>J+d$. If $\mathcal K_R=\prod_jW_j(x_j)F(R\prod_jx_j^{a_j})$, fixed $a_j$, and
\[
 \|F\|_{A,q}=\max_{h\le q}\sup_{u>0}(1+u)^A|(u\partial_u)^hF(u)|<\infty,
\]
then its Mellin coefficient has the pointwise bound
\[
 |b_R(\boldsymbol t)|\ll(1+R)^{-A}(1+|\boldsymbol t|)^{-q}
 \|F\|_{A,q}\prod_j\|W_j\|_{C^q}.
\]
\end{lemma}
\begin{proof}
Take the Fourier transform of $\mathcal K(e^{y_1},\ldots,e^{y_d})$ with factor $(2\pi)^{-d}$. Fourier inversion gives the identity. Applying $(1-\Delta_y)^{q/2}$ and integrating by parts gives the pointwise $q$-power decay. Ordinary and logarithmic derivative norms are comparable on the fixed box. In the special form the chain rule preserves $(1+R)^{-A}$, since the argument of $F$ is comparable to $R$. The weighted integral converges when $q>J+d$.
\end{proof}

\begin{lemma}[Row-dependent kernels; p2 Lemma B.2]\label{lem:rowkernel}
Let $I\Subset\operatorname{int}I_*$, and let $S_{j,r}(U)=\sum_na_{j,r}(n)U(x_{j,r,n})$, $j=1,2$, be arbitrary finite coefficient arrays, with $x_{j,r,n}>0$ and nonnegative row weights $w_r$. Suppose for every $U\in C_c^\infty(I_*)$ that
\[
 \sum_rw_r|S_{j,r}(U)|^2\le M_j\|U\|_{C^m(I_*)}^2.
\]
For $|c_r|\le w_r$ and kernels $\mathcal K_r$ supported in $I^2$,
\[
 \left|\sum_rc_r\sum_{n_1,n_2}a_{1,r}(n_1)\bar a_{2,r}(n_2)
 \mathcal K_r(x_{1,r,n_1},x_{2,r,n_2})\right|
 \ll_{I,I_*,m}\sqrt{M_1M_2}\sup_r\|\mathcal K_r\|_{C^{2m+4}(I^2)}.
\]
\end{lemma}
\begin{proof}
Choose a smooth $V$ supported in $I_*$ and equal to $1$ on $I$, and put $U_t(x)=V(x)x^{it}$. Mellin separation expresses each kernel as $\int b_r(s,t)U_s(x)\bar U_t(y)\dd s\dd t$, with common majorant
$\sup_r|b_r(s,t)|\ll\sup_r\|\mathcal K_r\|_{C^q}(1+|s|+|t|)^{-q}$. Take this common bound before summing rows and apply row Cauchy--Schwarz. The hypothesis bounds the two resulting norms by $\sqrt{M_j}(1+|t|)^m$. The remaining two-dimensional integral converges for $q=2m+4$. In particular the lemma does not allow choosing an arbitrary unrelated test separately for each row; the uniform derivative bound is what supplies that step.
\end{proof}

\begin{lemma}[Weighted recombination; p2 Lemma B.3]\label{lem:recombine}
For finite rows and labels, $w_r\ge0$, suppose $|c_{\iota,r}(\boldsymbol t)|\le b_\iota(\boldsymbol t)$ and $M=\sum_\iota\int b_\iota<\infty$. Then
\[
 \sum_rw_r\left|\sum_\iota\int c_{\iota,r}(\boldsymbol t)F_{\iota,\boldsymbol t}(r)\dd\boldsymbol t\right|^2
 \le M\sum_\iota\int b_\iota(\boldsymbol t)\sum_rw_r|F_{\iota,\boldsymbol t}(r)|^2\dd\boldsymbol t.
\]
The corresponding finite-sum assertion also holds.
\end{lemma}
\begin{proof}
Apply weighted Cauchy--Schwarz for each row with weights $b_\iota(\boldsymbol t)$, and then sum the nonnegative row weights. Tonelli applies to the nonnegative majorant when the right side is finite.
\end{proof}

\section{An explicit infinite-order character}\label{app:example}
The conductor-one character $\eta_6((z))=\alpha(z)^6$ is independent of the generator because every unit has sixth power $1$. At $\pi=-2-3\omega$, exact ring multiplication gives
\[
 \N(\pi)=7,\qquad\pi^3=19+18\omega,\qquad\pi^6=37+360\omega=-143+180i\sqrt3.
\]
Thus $\eta_6((\pi))=(-143+180i\sqrt3)/343$ has field trace $-286/343$ and is not an algebraic integer, hence not a root of unity. The character has infinite order. Its Euler factors at $7,5,3$ are, respectively,
\[
 (1+\tfrac{286}{343}7^{-s}+7^{-2s})^{-1},\qquad
 (1-5^{-2s})^{-1},\qquad(1+3^{-s})^{-1}.
\]
The algebraic character $z^6$ has arithmetic boundary $3+11/12=47/12$. These exact calculations identify the character, and Theorem~\ref{thm:A} gives its zero-free half-plane.

\bibliographystyle{plain}
\bibliography{references}
\section*{Use of AI}
The proofs and the text of this paper were produced with AI systems (GPT-6 Astra, GPT-6.1 Sol and Claude Opus~5.5) under the author's direction. They were checked by independent reviews with AI models and, for the explicit identities, by the exact and high-precision computations in the ancillary files. The author takes responsibility for the content.
\end{document}
